% !TeX root = ../../Book2.tex
% 纲要标题：## 第21章　Stiefel–Whitney、陈与 Pontryagin 类
% 纲要位置：几何学计划纲要.md，第 1802 行。
% 正文状态：初稿；尚未正式审核。

% 纲要内容开始
% 
% > 本章目标与证明范围：分别完成整数类与模2类的构造，再证明分裂原理和计算法。拓扑扭信息在此保留，不能用下一章的曲率形式取代。
% 
% 纲要内容结束

\chapter{Stiefel–Whitney、陈与 Pontryagin 类}
\label{g2:ch:21}
\BookChapterNumber{g2:ch:21}

\begin{chapterabstract}
	一张向量丛在不同维数上的扭转，可以由底空间上不同次数的上同调类记录。本章从实、复线丛出发，用射影丛的上同调关系构造 Stiefel--Whitney 类与陈类，再由复化定义 Pontryagin 类。Leray--Hirsch 定理和旗丛提供分裂原理，使线丛计算能够回到原来的向量丛。随后证明 Whitney 公式、顶类与 Euler 类的比较、Thom 平方解释和 Wu 公式，计算射影空间切丛与结构群提升障碍。最后在复 Grassmann 流形上把示性类与 Schubert 条件联系起来，完成四维复空间中二维子空间的上同调环及四条直线的相交计数。
\end{chapterabstract}

\begin{learninggoals}
	\begin{itemize}
		\item 将实线丛的方向输运、复线丛的 Euler 类和过渡函数联系起来，固定第一示性类的符号。
		\item 证明 Leray--Hirsch 定理、射影丛公式和分裂原理，说明形式根计算为何能决定原底空间上的等式。
		\item 构造 Stiefel--Whitney 类、陈类与 Pontryagin 类，掌握自然性、直和及复化公式，并辨认整数二阶扭信息。
		\item 完成 Thom 平方与 Wu 类的比较，用切丛计算检验向量场、几乎复结构和 spin 提升的条件。
		\item 在 Grassmann 流形上解释 Schubert 类及特殊类，计算 \(\operatorname{Gr}_2(\mathbb C^4)\) 的乘法并求解四条直线的相交问题。
	\end{itemize}
\end{learninggoals}

上一章的 Euler 类记录一张丛的截面能否避开零点。它的次数等于丛的实秩，因此只提供一个特定次数的信息：线丛的方向会沿闭路改变，较高秩丛又可能在二维或四维子空间上发生扭转，这些现象需要不同次数的类。示性类把一张丛组织成一组可拉回的上同调数据；同一个定义可以随映射限制到曲线、曲面或更高维的检验空间。

\ProseParagraphBreak

先考虑复秩二丛 \(E=L_1\oplus L_2\)。如果两条线的第一类分别为 \(x_1,x_2\)，它们的和与乘积就是最自然的对称数据：交换两条线不改变它们。但是一般丛并没有给定的线丛分解。我们会将 \(E\) 拉回到旗丛，在那里得到实际分解，并证明上同调拉回单射。于是在线丛上算出的等式能够返回原基空间。读者需要区分这两件事：分解发生在新的基空间，而等式由单射保证仍能检测原来的类。

\ProseParagraphBreak

为了避免先假定乘法公式再用它定义示性类，本章从射影丛的上同调关系读出各阶系数，再证明这些系数具有所需性质。实线丛使用模二一次类，复线丛使用整数二次类，两种构造具有相同的计算形状，系数却保留不同的信息。Pontryagin 类由实丛的复化得到；它的整数直和公式中还可能出现被二消去的修正项。我们将写出这个修正，并在实射影空间上给出非零例子。

\ProseParagraphBreak

本章使用第19章的丛分类、第20章的 Thom 与 Euler 构造，以及第13章的平方运算；所需条件会在具体证明中回引。分裂、平方与三类示性类的组织可参阅\cite[第23章第6--7节]{May1999}。最后一节的 Grassmann 计算把此前的上同调环、对偶和横截相交放在同一个空间里使用，也为后续代数几何中的交叉计算准备一个完整例子。

% !TeX root = ../../Book2.tex
% 纲要标题：### 21.1 线丛与通用生成元
% 纲要位置：几何学计划纲要.md，第 1806 行。
% 正文状态：未撰写；本文件目前是框架。

% 纲要内容开始
% 
% * 实线丛和第一 Stiefel–Whitney 类
% * 复线丛与第一 Chern 类
% * 过渡函数、分类空间和指数序列的比较
% 
% 纲要内容结束

\section{线丛与通用生成元}
\label{g2:sec:2101}

待撰写

\subsection{实线丛和第一 Stiefel–Whitney 类}
\label{g2:subsec:2101:01}
待撰写

\subsection{复线丛与第一 Chern 类}
\label{g2:subsec:2101:02}
待撰写

\subsection{过渡函数、分类空间和指数序列的比较}
\label{g2:subsec:2101:03}
待撰写
% 纲要条目排为编号小节标题，并进入目录与书签；说明段仍保留于纲要注释中。

% !TeX root = ../../Book2.tex
% 纲要标题：### 21.2 射影丛与分裂原理
% 纲要位置：几何学计划纲要.md，第 1812 行。
% 正文状态：未撰写；本文件目前是框架。

% 纲要内容开始
% 
% * Leray–Hirsch 的证明
% * 射影丛上同调
% * 旗丛拉回的单射性
% * 以形式根进行对称多项式计算
% 
% 纲要内容结束

\section{射影丛与分裂原理}
\label{g2:sec:2102}

待撰写

\subsection{Leray–Hirsch 的证明}
\label{g2:subsec:2102:01}
待撰写

\subsection{射影丛上同调}
\label{g2:subsec:2102:02}
待撰写

\subsection{旗丛拉回的单射性}
\label{g2:subsec:2102:03}
待撰写

\subsection{以形式根进行对称多项式计算}
\label{g2:subsec:2102:04}
待撰写
% 纲要条目排为编号小节标题，并进入目录与书签；说明段仍保留于纲要注释中。

% !TeX root = ../../Book2.tex
% 纲要标题：### 21.3 一般示性类的构造
% 纲要位置：几何学计划纲要.md，第 1819 行。
% 正文状态：初稿；尚未正式审核。

% 纲要内容开始
% 
% * 陈类的存在唯一性
% * Stiefel–Whitney 类及实射影丛关系
% * Whitney 和公式与 Wu 公式
% * 对偶、张量和复化
% * Pontryagin 类及其性质
% 
% 纲要内容结束

\section{一般示性类的构造}
\label{g2:sec:2103}

射影丛公式把一张秩 \(r\) 丛的粘合信息写入一个首一多项式关系，旗丛又把这个多项式变成线性因子的乘积。现在用关系的系数定义一般示性类，再利用旗丛拉回的单射性，把线丛上的计算送回原基空间。这一顺序同时解决存在性、选择无关性与计算问题。

以下主要以 CW 复形为基，有限秩丛均带纤维内积。证明中的射影丛及旗丛使用上一节的弱 CW 逼近比较。Thom 类也可沿同一比较运输：盘丛、球丛均为局部平凡丛，拉回到 CW 弱逼近后，总空间映射由纤维化同伦长正合列为弱等价；绝对上同调比较及对的长正合列给出相对上同调同构。因此第20章的 Thom、Euler 及其自然性可用于这些实际拉回丛。

\subsection{陈类的存在唯一性}
\label{g2:subsec:2103:01}
设 \(E\to B\) 为复秩 \(r\ge1\) 丛。在 \(p:\mathbb P(E)\to B\) 上记典范线丛为 \(\lambda\)，置 \(h=c_1(\lambda^*)\)。\cref{g2:thm:projective-bundle}说明 \(1,h,\ldots,h^{r-1}\) 是 \(H^*(B;\mathbb Z)\) 模基。因此 \(h^r\) 可以唯一写成这些幂的线性组合；其各项的总次数必须同为 \(2r\)。

\begin{definition}\label{g2:def:chern-classes-projective}
	复秩 \(r\) 丛 \(E\to B\) 的\term[chen-lei]{陈类}[Chern classes]为唯一满足
	\begin{equation}\label{g2:eq:chern-projective-relation}
		h^r+p^*c_1(E)h^{r-1}+p^*c_2(E)h^{r-2}+\cdots+p^*c_r(E)=0
	\end{equation}
	的类 \(c_i(E)\in H^{2i}(B;\mathbb Z)\)。置 \(c_0(E)=1\)，\(i>r\) 时置 \(c_i(E)=0\)，并称 \(c(E)=1+c_1(E)+\cdots+c_r(E)\) 为总陈类。
\end{definition}

当 \(r=1\) 时，\(\mathbb P(E)=B\)、\(\lambda=E\)，关系变成 \(-c_1(E)+c_1(E)=0\)，故新定义延续上一节的线丛规范。正号排列来自 \(h=c_1(\lambda^*)\)：如果改用 \(c_1(\lambda)=-h\)，关系的系数就要交错变号。

\begin{theorem}\label{g2:thm:characteristic-existence-uniqueness}
	CW 复形上的复有限秩向量丛具有上述自然陈类。若 \(q:Y\to B\) 是使 \(q^*E=L_1\oplus\cdots\oplus L_r\) 的旗丛，记 \(x_j=c_1(L_j)\)，则
	\begin{equation}\label{g2:eq:chern-root-formula}
		q^*c_i(E)=\sum_{j_1<\cdots<j_i}x_{j_1}\cdots x_{j_i},\qquad
		q^*c(E)=\prod_{j=1}^r(1+x_j).
	\end{equation}
	自然性及这个分裂公式唯一确定全部 \(c_i\)。
\end{theorem}
\begin{proof}
	射影丛模基已经保证存在唯一的关系系数。若 \(f:A\to B\)，则拉回方块中的 \(\mathbb P(f^*E)\to\mathbb P(E)\) 将 \(\lambda\) 拉回为新的典范线丛，故也将 \(h\) 拉回为新的 \(h\)。拉回式\eqref{g2:eq:chern-projective-relation}并用新射影丛的模基唯一性，得到 \(c_i(f^*E)=f^*c_i(E)\)。

	为确定关系的系数，先设 \(E=L_1\oplus\cdots\oplus L_r\) 已分裂。在 \(\mathbb P(E)\) 上，丛
	\[
		\operatorname{Hom}_{\mathbb C}(\lambda,p^*E)
		=\bigoplus_{j=1}^r(\lambda^*\otimes p^*L_j)
	\]
	有一个处处非零截面：每点的直线包含映射 \(\lambda_\ell\hookrightarrow E_{p(\ell)}\)。\cref{g2:prop:euler-basic-properties}说明它的底层定向实丛 Euler 类为零；同一命题的直和公式及线丛的 \(c_1=e_{\mathrm{Th}}\) 给出
	\[
		0=\prod_{j=1}^r c_1(\lambda^*\otimes p^*L_j)
		=\prod_{j=1}^r(h+p^*x_j).
	\]
	这里使用的是已证明的 Euler 类乘法，并未预先假定一般陈类的乘法。展开右端，与唯一的首一射影关系比较，得到 \(c_i(E)\) 是 \(x_j\) 的第 \(i\) 个初等对称式。

	一般情形拉回到旗丛，用自然性及刚才的分裂计算得到式\eqref{g2:eq:chern-root-formula}。由\cref{g2:thm:splitting-principle}，\(q^*\) 单射，所以这些等式唯一确定原基空间上的类；任何其他自然赋类若满足同样的分裂公式，拉回后相同，原来也必相同。
\end{proof}

例如秩二时，形式根 \(x_1,x_2\) 给出 \(c_1=x_1+x_2\)、\(c_2=x_1x_2\)，射影丛关系为 \(h^2+c_1h+c_2=0\)。根是旗空间上的实际一次线丛类，未必分别下降到底空间；能下降的这些对称式已经由射影关系独立构造。这是“形式根”计算可以可靠使用的依据。

\subsection{Stiefel–Whitney 类及实射影丛关系}
\label{g2:subsec:2103:02}
对实秩 \(r\) 丛，射影纤维是 \(\mathbb RP^{r-1}\)，用模二系数后它同样有一列由一次类的幂组成的基。这里不需要给原丛或射影纤维选整数方向。

\begin{definition}\label{g2:def:stiefel-whitney-projective}
	设 \(E\to B\) 为实秩 \(r\ge1\) 向量丛，\(p:\mathbb P(E)\to B\) 为实射影丛，\(\lambda\) 为典范实线丛，\(t=w_1(\lambda)\)。唯一满足
	\begin{equation}\label{g2:eq:sw-projective-relation}
		t^r+p^*w_1(E)t^{r-1}+p^*w_2(E)t^{r-2}+\cdots+p^*w_r(E)=0
	\end{equation}
	的类 \(w_i(E)\in H^i(B;\mathbb F_2)\)，称为 \(E\) 的\term[Stiefel-Whitney-lei]{Stiefel--Whitney 类}[Stiefel--Whitney classes]。置 \(w_0=1\)、\(w_i=0\)（\(i>r\)），并记 \(w(E)=\sum_iw_i(E)\)。
\end{definition}

\begin{theorem}\label{g2:thm:stiefel-whitney-construction}
	上述类对 CW 基空间上的实丛自然。若旗丛拉回使 \(q^*E=L_1\oplus\cdots\oplus L_r\)，并记 \(a_j=w_1(L_j)\)，则
	\begin{equation}\label{g2:eq:sw-root-formula}
		q^*w(E)=\prod_{j=1}^r(1+a_j).
	\end{equation}
	自然性及该分裂公式唯一确定这些类；秩一时与\cref{g2:def:first-stiefel-whitney-line}一致。
\end{theorem}
\begin{proof}
	射影丛的模二模基保证关系系数存在唯一，射影拉回方块保证自然性。若 \(E\) 已分裂，\(\operatorname{Hom}_{\mathbb R}(\lambda,p^*E)\) 仍有处处非零的包含截面。它分裂成 \(\lambda^*\otimes p^*L_j\) 的直和。由式\eqref{g2:eq:line-mod-two-euler}及模二 Euler 直和公式，得到
	\[
		0=e_2(\operatorname{Hom}(\lambda,p^*E))
		=\prod_{j=1}^r(t+p^*a_j).
	\]
	与式\eqref{g2:eq:sw-projective-relation}比较唯一系数，即得分裂公式。一般丛再拉回实旗丛；\(q^*\) 在模二上同调上单射，所以相同论证给出唯一性。秩一的关系为 \(t+w_1(E)=0\)，正好恢复线丛类。
\end{proof}

\begin{proposition}\label{g2:prop:first-classes-determinants}
	对 CW 基空间上的实丛 \(E\) 与复丛 \(F\)，有
	\[
		w_1(E)=w_1(\det E),\qquad c_1(F)=c_1(\det F).
	\]
	实丛可定向当且仅当 \(w_1(E)=0\)。
\end{proposition}
\begin{proof}
	在相应旗丛上，行列式线丛是各线因子的张量积，因此其第一类是各根之和；这也是两个总类的一次项。拉回单射使等式在底空间成立。实丛方向的输运为过渡矩阵的行列式符号，恰好是 \(\det E\) 的方向输运，所以可定向性由\cref{g2:prop:real-line-first-class}判定。
\end{proof}

\begin{theorem}\label{g2:thm:top-class-euler}
	若 \(E\to B\) 为 CW 基空间上的实秩 \(r\) 丛，则 \(w_r(E)=e_2(E)\)。若 \(E\) 整数可定向，则 \(w_r(E)=\rho_2 e_{\mathrm{Th}}(E)\)。若 \(F\to B\) 为复秩 \(r\) 丛，则
	\[
		c_r(F)=e_{\mathrm{Th}}(F_{\mathbb R}),
	\]
	右侧采用复自然定向。
\end{theorem}
\begin{proof}
	先在旗丛上分裂成线丛。实情形中，模二 Euler 类是各线的 \(w_1\) 之积，恰为式\eqref{g2:eq:sw-root-formula}的最高项。复情形中，实丛的直和顺序按各复线的正向实二平面排列，正是复自然定向；Euler 直和公式使最高类为 \(x_1\cdots x_r\)，恰为 \(c_r\)。再由拉回单射得到原空间等式。整数定向丛的 Thom 类约化模二后是规范模二 Thom 类，零截面拉回与系数约化交换，得到中间断言。
\end{proof}

\subsection{Whitney 和公式与 Wu 公式}
\label{g2:subsec:2103:03}
两张丛直和后，在线丛分裂模型中只是把两组根放在一起。因此直和公式来自一个简单的乘法，但要让它在原基空间成立，需要同时分裂两张丛并保留拉回单射。

\begin{theorem}[Whitney 和公式]\label{g2:thm:whitney-sum}
	设 \(B\) 为 CW 复形。对其上的复有限秩丛 \(E,F\)，以及实有限秩丛 \(V,W\)，分别有
	\begin{equation}\label{g2:eq:whitney-total-formulas}
		c(E\oplus F)=c(E)c(F),\qquad w(V\oplus W)=w(V)w(W).
	\end{equation}
	两组类都在添上平凡直和因子时不变。自然性、线丛归一化、秩以上的消失及式\eqref{g2:eq:whitney-total-formulas}也唯一刻画各组类。
\end{theorem}
\begin{proof}
	先取 \(E\) 的旗丛，再取拉回 \(F\) 的旗丛。复合拉回仍单射，且两丛同时分裂，设根分别为 \(x_i,y_j\)。直和的根就是它们的合并，所以
	\[
		q^*c(E\oplus F)=\prod_i(1+x_i)\prod_j(1+y_j)
		=q^*(c(E)c(F)).
	\]
	单射性给出复情形。实旗丛及一次模二根给出完全相同的乘法论证。平凡丛分裂成第一类均为零的线丛，其总类为 \(1\)，从而得到稳定性。任何另一组满足所列性质的类，在旗丛上由线丛归一化与直和公式确定为同一个根乘积，故由单射性与本构造一致。
\end{proof}

短正合列 \(0\to E\to G\to F\to0\) 也满足这个公式，因为纤维内积给出 \(G\cong E\oplus E^\perp\)，而 \(E^\perp\to F\) 是丛同构。这里需要的是拓扑分裂，若原丛带有全纯等额外结构，不要求分裂保持那些结构。

\ProseParagraphBreak

现在完成第13章留下的 Thom 平方比较。记实秩 \(r\) 丛的模二 Thom 类为 \(U\)，用第20章同构定义 \(w_i^{\mathrm{Sq}}\) 使
\[
	\operatorname{Sq}^iU=p^*w_i^{\mathrm{Sq}}(E)\smile U.
\]
\cref{g2:prop:thom-square-properties}已经用平方的自然性与 Cartan 公式证明这组类的自然性、直和公式和秩以上消失。第20章提供的 Thom 同构及纤维积直和公式使同一论证适用于本节的基空间，余下只需核对线丛归一化。

\begin{theorem}\label{g2:thm:thom-square-comparison}
	对 CW 复形上的实有限秩丛 \(E\)，射影丛构造的 Stiefel--Whitney 类满足
	\begin{equation}\label{g2:eq:thom-sw-identity-proved}
		\operatorname{Sq}(U)=p^*w(E)\smile U.
	\end{equation}
	因此 \(w_i^{\mathrm{Sq}}(E)=w_i(E)\)。
\end{theorem}
\begin{proof}
	设 \(L\) 为实线丛，\(U\) 是它的一次模二 Thom 类。平方的最高项性质给出 \(\operatorname{Sq}^1U=U\smile U\)。记 \(j^*\) 为忘去相对条件，则零截面与盘丛投影互为同伦逆，故
	\[
		j^*U=p^*e_2(L).
	\]
	相对杯积与绝对—相对杯积相容，所以
	\[
		U\smile U=j^*U\smile U=p^*e_2(L)\smile U.
	\]
	Thom 同构的单射性及式\eqref{g2:eq:line-mod-two-euler}说明 \(w_1^{\mathrm{Sq}}(L)=e_2(L)=w_1(L)\)。其余正次项由秩一消失。于是这组平方类满足刚才唯一性定理的全部条件，必与射影丛构造相同。
\end{proof}

这一步把作用于总空间相对类的平方运算，识别为基空间上的示性类。对流形的对角嵌入应用它，便能把切丛类与第13章已经定义的 Wu 类联系起来。

\begin{theorem}[Wu 公式]\label{g2:thm:wu-formula}
	设 \(M\) 为连通闭光滑 \(n\) 维流形，\(v(M)\) 为\cref{g2:def:wu-class}按模二对偶定义的总 Wu 类。则
	\[
		w(\mathrm TM)=\operatorname{Sq}(v(M)).
	\]
	这个等式不要求 \(M\) 整数可定向。
\end{theorem}
\begin{proof}
	全部系数取 \(\mathbb F_2\)。先说明证明的目标：把对角线的 Thom 类平方化，一边读出法丛的 \(w\)，另一边用对偶基读出 Wu 类的平方。记 \(H=H^*(M;\mathbb F_2)\)，并以 \(\epsilon(a)\) 表示 \(a\) 的 \(n\) 次分量在 \([M]_2\) 上的求值。闭流形有有限 CW 同伦型，故 \(H\) 有限维；Poincaré 对偶给出齐次基 \(b_\alpha\) 及其对偶基 \(b^\alpha\)，满足
	\[
		\epsilon(b_\alpha b^\beta)=\delta_\alpha^\beta.
	\]
	对角嵌入 \(\Delta:M\hookrightarrow M\times M\) 的法丛同构于 \(\mathrm TM\)：商线性映射 \((v,w)\mapsto w-v\) 的核就是对角切空间。由管状邻域的 Thom 类定义推前 \(\Delta_!a\)，即先在邻域内取 \(a\) 与法丛 Thom 类的乘积，再经切除和忘去相对条件得到 \(H^{n+\deg a}(M\times M;\mathbb F_2)\) 中的类。第20章法向局部对偶的同一计算给出
	\[
		\langle\Delta_!a\smile z,[M\times M]_2\rangle
		=\langle a\smile\Delta^*z,[M]_2\rangle.
	\]
	令 \(D=\Delta_!1\)。由域系数 Künneth 同构和上式的完美配对，
	\begin{equation}\label{g2:eq:diagonal-dual-basis}
		D=\sum_\alpha b_\alpha\times b^\alpha.
	\end{equation}
	例如将右边与 \(x\times y\) 相乘后求值，得到 \(\sum_\alpha\epsilon(b_\alpha x)\epsilon(b^\alpha y)=\epsilon(xy)\)，恰与对角拉回相同，故这确是 \(D\)。

	自然性使平方与相对类到绝对类的映射交换；\cref{g2:thm:thom-square-comparison}于是给出
	\[
		\operatorname{Sq}(D)=\Delta_!w(\mathrm TM)
		=D\smile\operatorname{pr}_1^*w(\mathrm TM).
	\]
	最后一个等式是邻域内的杯积公式：\(\operatorname{pr}_1\) 限制在管状邻域后，与收缩到对角线再投影同伦，所以其拉回正是 Thom 同构中的基空间因子。

	在 \(H\otimes H\) 上施加线性映射 \(1\otimes\epsilon\)。对右边，式\eqref{g2:eq:diagonal-dual-basis}给出 \((1\otimes\epsilon)D=1\)，所以结果是 \(w(\mathrm TM)\)。对左边，Cartan 公式和外积相容性给出
	\[
		(1\otimes\epsilon)\operatorname{Sq}(D)
		=\sum_\alpha\operatorname{Sq}(b_\alpha)\,\epsilon(\operatorname{Sq}(b^\alpha)).
	\]
	Wu 类的定义逐次数合起来正是 \(\epsilon(\operatorname{Sq}(x))=\epsilon(v(M)x)\)。因此上式等于
	\[
		\operatorname{Sq}\left(\sum_\alpha b_\alpha\,\epsilon(v(M)b^\alpha)\right)
		=\operatorname{Sq}(v(M)),
	\]
	最后一步是对偶基展开。两边比较即得所求公式，全程只使用模二基本类。
\end{proof}

这就证明了\cref{g2:thm:wu-tangent-preview}的预告。低次时，\(w_1=v_1\)、\(w_2=v_2+\operatorname{Sq}^1v_1\)；后者额外的一项说明 Wu 类本身与切丛类一般不同。第13章中 \(\mathbb RP^2\) 与 \(\mathbb CP^2\) 的计算现在分别给出实际切丛的总类，而不再只是运算产生的候选表达。

\subsection{对偶、张量和复化}
\label{g2:subsec:2103:04}
对偶、张量与复化改变线丛的根，因此可以用同一个分裂过程计算。进行比较时须保留系数：复陈类位于整数上同调，实 Stiefel--Whitney 类位于模二上同调，二者之间使用约化映射 \(\rho_2\)。

\begin{proposition}\label{g2:prop:chern-operations}
	设 \(E,F\) 为 CW 基空间上的复丛，秩分别为 \(r,s\)，形式陈根为 \(x_i,y_j\)。则
	\[
		c_k(E^*)=c_k(\overline E)=(-1)^k c_k(E),\qquad
		c(E\otimes F)=\prod_{i,j}(1+x_i+y_j),
	\]
	第二式在同时分裂空间中计算，再由单射性理解于底空间。特别地，
	\[
		c_1(E\otimes F)=s\,c_1(E)+r\,c_1(F).
	\]
	若 \(L\) 为复线丛、\(\ell=c_1(L)\)，则
	\begin{equation}\label{g2:eq:chern-line-twist}
		c_k(E\otimes L)=\sum_{j=0}^k\binom{r-j}{k-j}c_j(E)\ell^{k-j}.
	\end{equation}
\end{proposition}
\begin{proof}
	在同时分裂空间，\(E^*\)、\(\overline E\) 的线因子第一类为 \(-x_i\)，因为线丛对偶与共轭都变号；第 \(k\) 个初等对称式遂乘以 \((-1)^k\)。张量积实际分裂为全部 \(L_i\otimes K_j\)，每个因子的第一类为 \(x_i+y_j\)，给出乘积公式及其一次项。

	对线丛扭转，把第 \(k\) 个初等对称式写为 \(e_k(x_1+\ell,\ldots,x_r+\ell)\)。从一项的 \(k\) 个因子中选 \(j\) 个贡献 \(x\)，其余贡献 \(\ell\)；固定那 \(j\) 个根后，还可从其余 \(r-j\) 个位置选出 \(k-j\) 个，故系数为 \(\binom{r-j}{k-j}\)。这证明式\eqref{g2:eq:chern-line-twist}。各等式都是对称式，且拉回单射，因此在底空间成立。
\end{proof}

实丛的对偶由内积与自身同构，所以 \(w(E^*)=w(E)\)。若两实丛分别分裂为根 \(a_i,b_j\) 的实线丛，则其张量积满足 \(w(E\otimes F)=\prod_{i,j}(1+a_i+b_j)\)。这些公式的乘积次数不同，不能把总陈类的整数系数式直接改名为总 Stiefel--Whitney 类的式子。

\begin{proposition}\label{g2:prop:real-complex-characteristic}
	对 CW 基空间上的复秩 \(r\) 丛 \(F\)，有
	\begin{equation}\label{g2:eq:complex-real-sw}
		w_{2i}(F_{\mathbb R})=\rho_2c_i(F),\qquad w_{2i+1}(F_{\mathbb R})=0.
	\end{equation}
	对实丛 \(E\) 的复化 \(E_{\mathbb C}=E\otimes_{\mathbb R}\mathbb C\)，有
	\begin{equation}\label{g2:eq:complexification-chern-mod-two}
		\rho_2c_i(E_{\mathbb C})=w_i(E)^2,\qquad 2c_{2j+1}(E_{\mathbb C})=0.
	\end{equation}
\end{proposition}
\begin{proof}
	在复旗丛上分裂 \(F\)。该拉回不仅在整数上同调单射，在模二上同调也单射，因为复射影纤维的幂基约化后仍是一组基。每条复线的底层实二平面定向，故 \(w_1=0\)，而\cref{g2:thm:top-class-euler}给出 \(w_2=\rho_2c_1\)。对各线因子应用 Whitney 公式，得 \(w(F_{\mathbb R})=\prod_i(1+\rho_2x_i)\)，比较次数即得第一组公式。

	映射 \(v+\mathrm i w\mapsto(v,w)\) 给出实丛同构 \((E_{\mathbb C})_{\mathbb R}\cong E\oplus E\)，于是
	\[
		\sum_i\rho_2c_i(E_{\mathbb C})=w(E)^2=\sum_iw_i(E)^2
	\]
	在模二系数下成立；最后一步的混合项成对相消。比较次数得到第二组的第一个式子。复化丛有规范的复线性同构 \(E_{\mathbb C}\cong\overline{E_{\mathbb C}}\)，由共轭的陈公式得 \(c_i=(-1)^ic_i\)，故奇下标项被 \(2\) 消去。
\end{proof}

例如实线丛复化后虽然只有一个陈类，这个类也未必为零：式\eqref{g2:eq:line-complexification-bockstein}精确把它写成方向类的整数 Bockstein。其模二约化为 \(\operatorname{Sq}^1w_1=w_1^2\)，与式\eqref{g2:eq:complexification-chern-mod-two}一致。复化保留了某些实方向信息，但把它从一次模二群移入二次整数扭群。

\subsection{Pontryagin 类及其性质}
\label{g2:subsec:2103:05}
实丛复化以后得到整数陈类。奇下标类可能只剩二阶扭信息，偶下标类则给出位于 \(4i\) 次的 Pontryagin 类。定义中的交错符号使定向实二平面的第一 Pontryagin 类成为 Euler 类的平方。

\begin{definition}\label{g2:def:pontryagin-classes}
	实有限秩丛 \(E\to B\) 的\term[Pontryagin-lei]{Pontryagin 类}[Pontryagin classes]定义为
	\begin{equation}\label{g2:eq:pontryagin-definition}
		p_i(E)=(-1)^i c_{2i}(E_{\mathbb C})\in H^{4i}(B;\mathbb Z),\qquad p(E)=1+p_1(E)+p_2(E)+\cdots.
	\end{equation}
	若 \(2i>\operatorname{rank}_{\mathbb R}E\)，则该项为零。
\end{definition}

Pontryagin\footnote{\OriginalName{Лев Семёнович Понтрягин}（Lev Semenovich Pontryagin，1908--1988），苏联数学家，研究代数拓扑、拓扑群与最优控制理论。}类对实丛定义，不要求它可定向。

\begin{proposition}\label{g2:prop:pontryagin-properties}
	对 CW 基空间上的实丛，Pontryagin 类自然且在加平凡直和因子后不变，并满足
	\[
		\rho_2p_i(E)=w_{2i}(E)^2.
	\]
	若 \(E\) 是定向实秩 \(2r\) 丛，则 \(p_r(E)=e_{\mathrm{Th}}(E)^2\)。若 \(F\) 为复丛，形式陈根为 \(x_j\)，则
	\begin{equation}\label{g2:eq:pontryagin-complex-underlying}
		p(F_{\mathbb R})=\prod_j(1+x_j^2),\qquad
		p_1(F_{\mathbb R})=c_1(F)^2-2c_2(F).
	\end{equation}
\end{proposition}
\begin{proof}
	自然性、稳定性来自复化与陈类的相应性质。模二约化时符号消失，再用式\eqref{g2:eq:complexification-chern-mod-two}即可。

	对定向秩 \(2r\) 丛 \(E\)，实同构 \((E_{\mathbb C})_{\mathbb R}\cong E\oplus E\) 所比较的两种方向须明确。若 \(e_1,\ldots,e_{2r}\) 为正向实基，复自然实方向是交错排列 \((e_1,\mathrm i e_1,\ldots,e_{2r},\mathrm i e_{2r})\)，直和方向则先放全部 \(e_j\)，再放全部 \(\mathrm i e_j\)。它们相差 \((-1)^{(2r)(2r-1)/2}=(-1)^r\)。最高陈类等于复方向 Euler 类，因而
	\[
		c_{2r}(E_{\mathbb C})=(-1)^re_{\mathrm{Th}}(E\oplus E)=(-1)^re_{\mathrm{Th}}(E)^2.
	\]
	代入定义便得 \(p_r=e_{\mathrm{Th}}^2\)。

	最后，复线性同构 \((F_{\mathbb R})_{\mathbb C}\cong F\oplus\overline F\) 可在每根纤维由复结构的 \(\pm\mathrm i\) 特征子空间给出。因此其陈总类在分裂空间为
	\[
		\prod_j(1+x_j)(1-x_j)=\prod_j(1-x_j^2).
	\]
	取偶下标项再乘 \((-1)^i\)，得到式\eqref{g2:eq:pontryagin-complex-underlying}。一次项利用 \(\sum_jx_j^2=(\sum_jx_j)^2-2\sum_{i<j}x_ix_j\) 化成所列公式。
\end{proof}

整数 Pontryagin 类的直和关系还要保留奇陈类的乘积。直接从陈类的 Whitney 公式分开偶、奇下标，可得
\begin{equation}\label{g2:eq:pontryagin-sum-torsion}
	\begin{split}
		p_k(E\oplus F)={}&\sum_{i+j=k}p_i(E)p_j(F)\\
		&+(-1)^k\!\sum_{\substack{a+b=2k\\a,b\text{ 为奇数}}}c_a(E_{\mathbb C})c_b(F_{\mathbb C}).
	\end{split}
\end{equation}
第二行被 \(2\) 消去，但未必等于零。因此
\[
	p(E\oplus F)=p(E)p(F)\quad\text{于 }H^*(B;\mathbb Z[1/2])
\]
总是成立；在整数上同调中，若相关群没有二阶扭，或者相关奇陈乘积全部消失，也可使用同一公式。特别是两个丛都来自复丛的底层实丛时，上面根的成对表达使这些奇项确实为零。

\begin{example}\label{g2:exm:pontryagin-whitney-torsion}
	令 \(L\to\mathbb RP^4\) 为典范实线丛，\(a=w_1(L)\)，\(t=c_1(L_{\mathbb C})=\widetilde\beta_2a\)。由第13章 \(\rho_2\widetilde\beta_2=\operatorname{Sq}^1\) 及最高平方性质，\(\rho_2t=a^2\)，所以 \(\rho_2(t^2)=a^4\ne0\)，从而 \(t^2\ne0\)。每条实线丛的 Pontryagin 类在正次全为零，但
	\[
		p_1(L\oplus L)=-c_2(L_{\mathbb C}\oplus L_{\mathbb C})=-t^2\ne0.
	\]
	这里整个差正是式\eqref{g2:eq:pontryagin-sum-torsion}的奇下标乘积。换成实或有理系数后 \(t\) 消失，便看不到这个例子中的区别。
\end{example}

至此，三个系列的次数与系数可以并列为
\[
\begin{array}{c|c|c|c}
	\text{丛的类型}&\text{类}&\text{所在群}&\text{构造所用数据}\\\hline
	\text{实秩 }r&w_i&H^i(B;\mathbb F_2)&\text{实射影丛关系}\\
	\text{复秩 }r&c_i&H^{2i}(B;\mathbb Z)&\text{复射影丛关系}\\
	\text{实秩 }r&p_i&H^{4i}(B;\mathbb Z)&(-1)^ic_{2i}(E_{\mathbb C})
\end{array}
\]
下一节将这些类作用于实际切丛和结构群约化；下一章再在实、复系数中用曲率给出代表。整数类及其扭项已经在拓扑构造中确定，曲率比较须与这里的规范一致。
% 纲要条目排为编号小节标题，并进入目录与书签；说明段仍保留于纲要注释中。

% !TeX root = ../../Book2.tex
% 纲要标题：### 21.4 几何计算与约束
% 纲要位置：几何学计划纲要.md，第 1827 行。
% 正文状态：未撰写；本文件目前是框架。

% 纲要内容开始
% 
% * 射影空间切丛
% * 球丛和低维切丛
% * 特征数与结构群约化
% * spin 条件与更高障碍的区别
% 
% 纲要内容结束

\section{几何计算与约束}
\label{g2:sec:2104}

待撰写

\subsection{射影空间切丛}
\label{g2:subsec:2104:01}
待撰写

\subsection{球丛和低维切丛}
\label{g2:subsec:2104:02}
待撰写

\subsection{特征数与结构群约化}
\label{g2:subsec:2104:03}
待撰写

\subsection{spin 条件与更高障碍的区别}
\label{g2:subsec:2104:04}
待撰写
% 纲要条目排为编号小节标题，并进入目录与书签；说明段仍保留于纲要注释中。

% !TeX root = ../../Book2.tex
% 纲要标题：### 21.5 Grassmann 流形的胞腔、Schubert 类与相交计算
% 纲要位置：几何学计划纲要.md，第 1834 行。
% 正文状态：初稿；尚未正式审核。

% 纲要内容开始
% 
% * 固定完全旗与 Schubert 胞腔的秩条件
% * 行阶梯形坐标、偶维胞腔与整数上同调基
% * 通用子丛、商丛和特殊 Schubert 类
% * 利用旗的移动及相交配对计算低阶乘法
% * $\operatorname{Gr}(2,4)$ 的杯积计算与四条直线问题
% 
%   在复数域模型上计算特殊除子类的四次幂，并由一般位置保证几何计数；完整的方程消去、重数与退化核验放在 V.27.5。本节不以一般 Schubert 簇光滑为前提，也不追加完整旗簇表示论。
% 
% * 与 V.27.5 中 Chow 环及 Plücker 方程计算的比较
% 
% ---
% 
% 纲要内容结束

\section{Grassmann 流形的胞腔、Schubert 类与相交计算}
\label{g2:sec:2105}

示性类已经能把向量丛的运算变成上同调计算。现在反过来问：这些上同调类是否有具体的几何代表，它们的杯积又在数什么？复 Grassmann 流形提供了一个可以同时看见两面的模型。一个子空间与固定旗相交的维数条件给出 Schubert 簇；通用商丛的陈类恰好代表其中最简单的一类条件。最后，\(\operatorname{Gr}(2,4)\) 上一个四次杯积将回答三维射影空间中四条直线的公共截线问题。

以下记 \(G=\operatorname{Gr}(k,n)=\operatorname{Gr}_k(\mathbb C^n)\)，其中 \(1\le k<n\)，并置 \(q=n-k\)、\(N=kq\)。它是复维数 \(N\) 的紧光滑流形，始终采用复结构给出的实定向及整数上同调。第1章的\cref{g2:thm:grassmann-schubert-cw}已经证明胞腔的特征映射及闭包关系；这里恢复坐标含义，再用它们构造同调类并计算乘法。有关胞腔、相交与通用丛的经典叙述，可参阅\cite[第1章第5节，The Cell Decomposition、The Schubert Calculus、Universal Bundles]{GriffithsHarris1994}。

\subsection{固定完全旗与 Schubert 胞腔的秩条件}
\label{g2:subsec:2105:01}
固定一条\term[wanquanqi]{完全旗}[complete flag]
\[
	0=F_0\subset F_1\subset\cdots\subset F_n=\mathbb C^n,
	\qquad\dim_{\mathbb C}F_i=i.
\]
选取与它相容的基 \(e_1,\ldots,e_n\)，使 \(F_i=\operatorname{span}(e_1,\ldots,e_i)\)。对 \(V\in G\)，整数 \(\dim(V\cap F_i)\) 从零增到 \(k\)，每步最多增加一，所以恰有 \(k\) 个跳跃位置
\[
	1\le d_1<d_2<\cdots<d_k\le n.
\]
第1章用这些递增指标编号胞腔。为了让指标直接记录余维，现在改用 \(\lambda_i=q+i-d_i\)。于是 \(q\ge\lambda_1\ge\cdots\ge\lambda_k\ge0\)。这样的整数列称为放在 \(k\times q\) 矩形中的一个分拆，记其大小为 \(|\lambda|=\sum_i\lambda_i\)。它可画成左对齐的方格行，行长依次为 \(\lambda_1,\ldots,\lambda_k\)；方格总数将成为复余维。

\begin{definition}\label{g2:def:schubert-partition-locus}
	给定上述旗与分拆，令 \(d_i=q+i-\lambda_i\)。具有跳跃位置恰为 \(d_1,\ldots,d_k\) 的子空间组成 Schubert 胞腔 \(C_\lambda(F_\bullet)\)。其闭包
	\begin{equation}\label{g2:eq:schubert-rank-conditions}
		\Omega_\lambda(F_\bullet)=\{V\in G:\dim(V\cap F_{d_i})\ge i,\ 1\le i\le k\}
	\end{equation}
	称为相对于该旗的\term[Schubert-cu]{Schubert 簇}[Schubert variety]。
\end{definition}

只写 \(\dim(V\cap F_{d_i})=i\) 还不足以定义开胞腔：还要排除跳跃提前发生。因此等价的精确条件是同时要求
\[
	\dim(V\cap F_{d_i})=i,\qquad\dim(V\cap F_{d_i-1})=i-1.
\]
例如 \(k=1\) 时，\(C_{(j)}=\mathbb P(F_{n-j})\setminus\mathbb P(F_{n-j-1})\)，而 \(\Omega_{(j)}=\mathbb P(F_{n-j})\) 是复射影空间中的余维 \(j\) 线性子空间。闭包中的不等式允许跳跃提前发生；它们比开胞腔的等式更适合描述极限。

\ProseParagraphBreak

这些条件也确实是矩阵的秩条件。若 \(V\) 由一个满秩 \(n\times k\) 矩阵 \(A\) 的列张成，则
\[
	\dim(V\cap F_d)=k-\operatorname{rank}(\mathbb C^k\xrightarrow{A}\mathbb C^n\longrightarrow\mathbb C^n/F_d).
\]
因而 \(\dim(V\cap F_d)\ge i\) 等价于最后 \(n-d\) 行组成的矩阵秩不超过 \(k-i\)，可由全部 \((k-i+1)\) 阶子式为零表达。这里条件与所选 \(V\) 的基无关，因为更换基只在右侧乘一个可逆矩阵。它说明 Schubert 簇在 Grassmann 坐标中由多项式方程给出，也说明闭集中的奇点为何可能出现：一个秩条件的不同阶退化会在同一闭包中相遇。

\subsection{行阶梯形坐标、偶维胞腔与整数上同调基}
\label{g2:subsec:2105:02}
在 \(C_\lambda\) 内，从右向左消元得到唯一的基
\begin{equation}\label{g2:eq:schubert-echelon-coordinates}
	v_i=e_{d_i}+\sum_{\substack{a<d_i\\a\notin\{d_1,\ldots,d_{i-1}\}}}z_{ia}e_a,
	\qquad 1\le i\le k.
\end{equation}
第 \(i\) 个向量有 \(d_i-i=q-\lambda_i\) 个自由复系数，因而
\[
	C_\lambda\cong\mathbb C^{N-|\lambda|},\qquad
	\Omega_\lambda=\coprod_{\mu_i\ge\lambda_i\ \text{对所有 }i}C_\mu.
\]
胞腔的实维数是 \(2(N-|\lambda|)\)。例如在 \(\operatorname{Gr}(2,4)\) 中，\(\lambda=(1,0)\) 对应 \((d_1,d_2)=(2,4)\)，可以把一组基作为矩阵的行写成
\[
	\begin{pmatrix}z_{11}&1&0&0\\z_{21}&0&z_{23}&1\end{pmatrix}.
\]
三个自由复数给出六维实胞腔；分拆只有一个方格，故它在八维实 Grassmann 流形中余维二。

\begin{proposition}\label{g2:prop:schubert-integral-basis}
	复 Grassmann 流形 \(G=\operatorname{Gr}(k,n)\) 的整数同调与上同调没有扭元，奇数次群均为零。每个 \(\Omega_\lambda\) 有一个由其开胞腔复定向确定的类
	\[
		[\Omega_\lambda]\in H_{2(N-|\lambda|)}(\Omega_\lambda;\mathbb Z).
	\]
	把这些类推入 \(G\)，得到各次同调的基；其 Poincar\'e 对偶
	\begin{equation}\label{g2:eq:schubert-class-definition}
		\sigma_\lambda=\mathrm{PD}_G[\Omega_\lambda]\in H^{2|\lambda|}(G;\mathbb Z)
	\end{equation}
	称为\term[Schubert-lei]{Schubert 类}[Schubert class]，并构成整数上同调的加法基。这些类与所选完全旗无关。
\end{proposition}
\begin{proof}
	由\cref{g2:thm:grassmann-schubert-cw}，上述集合给出有限 CW 结构，且全部胞腔的实维数为偶数。胞腔微分降低一次，所以其源或靶必为零，全部微分为零。于是同调在每个偶数维自由生成于该维胞腔，在奇数维为零。万有系数定理中的 Ext 项也为零，得到对应的上同调结论。

	\(\Omega_\lambda\) 是一个子复形，唯一的最高维胞腔为 \(C_\lambda\)，其余胞腔的实维数至少低二。故其最高维同调为 \(\mathbb Z\)，到最高维胞腔相对同调的映射为同构。选择在 \(C_\lambda\) 上给出复定向生成元的唯一类，便定义了 \([\Omega_\lambda]\)。包含诱导的胞腔链映射将这个生成元送到 \(G\) 中同一胞腔的生成元，所以所有这类类组成同调基。这里的构造只用胞腔复形，并未要求整个闭包是流形。

	\(G\) 是闭定向流形，\cref{g2:cor:closed-poincare-duality}将这个同调基送到上同调基。任意两条完全旗由某个 \(g\in\operatorname{GL}_n(\mathbb C)\) 联系。该群道路连通：极分解先把正定因子收缩到恒等，酉因子再逐个缩小其特征值的辐角。因此可取从恒等到 \(g\) 的道路，给出 \(G\) 上恒等映射与 \(V\mapsto gV\) 的同伦。后一个映射保持各开胞腔的复定向，并把对应的 Schubert 簇相互送到，所以它们给出相同同调类及对偶类。
\end{proof}

几何相交时，还可以用一个实际的闭光滑流形来代表这个类。令 \(\widetilde\Omega_\lambda\) 由所有旗
\[
	0=W_0\subset W_1\subset\cdots\subset W_k,
	\qquad\dim W_i=i,\quad W_i\subset F_{d_i}
\]
组成，并令 \(\rho_\lambda(W_\bullet)=W_k\)。已选 \(W_{i-1}\) 后，选择 \(W_i\) 等价于在 \(F_{d_i}/W_{i-1}\) 中选一条线，因而每一步都是复射影丛，纤维为 \(\mathbb{CP}^{d_i-i}\)。这说明 \(\widetilde\Omega_\lambda\) 是紧光滑复流形，复维数为 \(\sum_i(d_i-i)=N-|\lambda|\)。

\ProseParagraphBreak

映射 \(\rho_\lambda\) 的像恰为 \(\Omega_\lambda\)：给定满足秩条件的 \(V\)，可在嵌套的 \(V\cap F_{d_i}\) 中逐次选择 \(W_i\)。若 \(V\in C_\lambda\)，则维数迫使 \(W_i=V\cap F_{d_i}\)，选择唯一；在这个开集上，常秩核空间随矩阵坐标光滑且复解析地变化。因此 \(\rho_\lambda\) 在开胞腔上是保持复定向的同构，而
\begin{equation}\label{g2:eq:schubert-resolution-class}
	(\rho_\lambda)_*[\widetilde\Omega_\lambda]=[\Omega_\lambda].
\end{equation}
最后一个等式可限制到最高维胞腔核对：两边都给出其正生成元，上一命题的唯一性使它们相等。这个旗模型将闭包低维处的多重选择留在纤维中，同时为相交理论提供一个光滑源。

\subsection{通用子丛、商丛和特殊 Schubert 类}
\label{g2:subsec:2105:03}
在 \(G\) 上有通用子丛 \(S\) 与商丛 \(Q\)：它们在 \(V\) 处的纤维分别为 \(V\) 和 \(\mathbb C^n/V\)。这是第19章典范丛的有限维模型，构成短正合列
\begin{equation}\label{g2:eq:grassmann-tautological-sequence}
	0\longrightarrow S\longrightarrow\underline{\mathbb C}^{\,n}\longrightarrow Q\longrightarrow0.
\end{equation}
选择 Hermitian 内积后可把 \(Q\) 识别为 \(S^\perp\)，所以 Whitney 和公式给出 \(c(S)c(Q)=1\)。这个等式已经把两张丛的陈类联系起来；以下识别它们所记录的几何条件。

\begin{theorem}\label{g2:thm:special-schubert-chern}
	在 \(G=\operatorname{Gr}(k,n)\) 上，令 \(q=n-k\)。对 \(1\le j\le q\)，有
	\begin{equation}\label{g2:eq:special-schubert-quotient}
		\sigma_j:=\sigma_{(j,0,\ldots,0)}=c_j(Q).
	\end{equation}
	特别地，\(\sigma_1=c_1(Q)=c_1(S^*)=c_1(\det S^*)\)。
\end{theorem}
\begin{proof}
	取 \(A=F_{q+1-j}\)，记 \(a=\dim A=q+1-j\)。对于分拆 \((j,0,\ldots,0)\)，除第一个条件外，其余条件由维数不等式自动成立。因此 \(\Omega_j=\{V:V\cap A\ne0\}\)。为把这个非零交条件写成一个横截零集，在 \(G\times\mathbb P(A)\) 上记 \(\ell\) 为从第二因子拉回的重言线丛，考虑秩 \(q\) 丛 \(\ell^*\otimes Q\) 的截面
	\[
		s(V,L):L\hookrightarrow\mathbb C^n\longrightarrow\mathbb C^n/V.
	\]
	其零集为 \(Z=\{(V,L):L\subset A\cap V\}\)。固定 \(L\) 后，\(V/L\) 可任取 \(\mathbb C^n/L\) 中的 \((k-1)\) 维子空间，所以 \(Z\) 是 \(\mathbb P(A)\) 上的光滑 Grassmann 丛。

	还可直接核对截面横截：在 \(L\subset V\) 处，\(G\) 的切空间由图坐标识别为 \(\operatorname{Hom}_{\mathbb C}(V,\mathbb C^n/V)\)。截面在这些方向的微分，除一个无关紧要的负号外，就是限制到 \(L\) 的映射
	\[
		\operatorname{Hom}(V,\mathbb C^n/V)\longrightarrow\operatorname{Hom}(L,\mathbb C^n/V),
	\]
	它满射，因为 \(L\) 上的线性映射可沿任意补空间延拓。故 \(s\) 横截于零截面。由\cref{g2:thm:top-class-euler}以及\cref{g2:thm:euler-transverse-zero-duality}，\(Z\) 的对偶类是 \(c_q(\ell^*\otimes Q)\)，定向就是复定向。

	令 \(h=c_1(\ell^*)\)。分裂原理把 \(Q\) 的形式根记为 \(x_1,\ldots,x_q\)，从而
	\begin{equation}\label{g2:eq:schubert-incidence-top-chern}
		c_q(\ell^*\otimes Q)=\prod_{b=1}^q(h+x_b)
		=\sum_{i=0}^q h^{q-i}c_i(Q).
	\end{equation}
	这是在旗丛拉回后验证的对称多项式恒等式；拉回单射使它在原空间成立。沿 \(\mathbb P(A)=\mathbb{CP}^{a-1}\) 求值时，只有 \(h^{a-1}=h^{q-j}\) 留下系数一，所以式\eqref{g2:eq:schubert-incidence-top-chern}给出 \(c_j(Q)\)。

	把这个系数提取写成同调推前，便可直接使用已证的乘积求值公式。记投影为 \(\pi:G\times\mathbb P(A)\to G\)。对任意 \(\beta\in H^{2(N-j)}(G;\mathbb Z)\)，顶陈类与零点类的比较给出
	\[
	\begin{aligned}
		\langle\beta,\pi_*[Z]\rangle
		&=\langle\pi^*\beta,[Z]\rangle\\
		&=\left\langle\pi^*\beta\smile\sum_i h^{q-i}c_i(Q),[G\times\mathbb{CP}^{a-1}]\right\rangle\\
		&=\langle\beta\smile c_j(Q),[G]\rangle.
	\end{aligned}
	\]
	全部次数均为偶数，交换因子没有符号。另一方面，\(\pi|_Z\) 的像为 \(\Omega_j\)，在其开胞腔上 \(V\cap A\) 恰为一维，故该处投影是保持复定向的同构。最高维胞腔的归一化给出 \(\pi_*[Z]=[\Omega_j]\)。\(G\) 的整数上同调自由，Poincar\'e 配对非退化，因此上述全部求值确定 \(\sigma_j=c_j(Q)\)。最后对式\eqref{g2:eq:grassmann-tautological-sequence}取一次陈类，并用对偶与行列式的公式，得到 \(\sigma_1\) 的其余表达。
\end{proof}

这个证明的几何步骤是先把“至少有一条交线”改成“连这条线也一起选出”。增加 \(\mathbb P(A)\) 参数后，退化条件成为一张向量丛的横截截面零集；随后沿这些参数求值，消去辅助选择，得到所求陈类。

\subsection{利用旗的移动及相交配对计算低阶乘法}
\label{g2:subsec:2105:04}
同一个 Schubert 类可以由许多不同旗给出的簇代表。移动旗的用途，是让相交发生在光滑开胞腔中并且横截；这时每个复交点贡献 \(+1\)，杯积便成为实际计数。

\begin{proposition}\label{g2:prop:schubert-moving-flags}
	设分拆 \(\lambda^{(1)},\ldots,\lambda^{(m)}\) 满足 \(\sum_b|\lambda^{(b)}|=N\)。可以选择处于一般位置的完全旗 \(F_\bullet^{(b)}\)，使各 \(\Omega_{\lambda^{(b)}}(F_\bullet^{(b)})\) 的公共交集是有限个点，全部位于各自开胞腔中且横截，并且
	\[
		\#\bigcap_{b=1}^m\Omega_{\lambda^{(b)}}(F_\bullet^{(b)})
		=\left\langle\prod_{b=1}^m\sigma_{\lambda^{(b)}},[G]\right\rangle.
	\]
	这里一般位置可取为旗参数空间中的一个开稠密集合；交集允许为空。
\end{proposition}
\begin{proof}
	群 \(\operatorname{GL}_n(\mathbb C)\) 在 \(G\) 上传递作用。作用映射 \((g,V)\mapsto gV\) 是淹没：在图坐标 \(\mathrm T_VG=\operatorname{Hom}(V,\mathbb C^n/V)\) 中，矩阵的一阶变化可以任意指定这张线性映射。因此将一个固定光滑胞腔随 \(g\) 移动，\cref{g2:thm:parametric-transversality}保证对一般参数，它横截于任意另一个固定光滑子流形。

	每个 Schubert 簇只有有限个胞腔。依次移动这些旗，同时对全部胞腔组合应用参数横截性，就使任意组合的相交都横截。若所有胞腔都取开胞腔，预期复维数为 \(N-\sum_b|\lambda^{(b)}|=0\)。若至少一个取闭包中的低维胞腔，其复维数至少再低一，预期维数变成负数，横截性迫使交集为空。所以全部交点落在开胞腔中。簇的闭性与 \(G\) 的紧性使总交集紧，而零维横截交集离散，故只有有限点。

	用式\eqref{g2:eq:schubert-resolution-class}的光滑旗模型代表各同调类。其乘积映射到 \(G^m\)，公共交点对应于这个映射对对角子流形的逆像；由于交点都在开胞腔中，那里各旗模型与 Schubert 簇相同，逆像仍然横截。将\cref{g2:thm:geometric-intersection-pairing}用于对角子流形，再用帽积的自然性及各个基本类的推前，所得求值就是所列杯积在 \([G]\) 上的值。这也把带奇点闭包的计数归结为闭光滑源与光滑对角的相交。

	在交点附近，所有切空间和法向商空间都是复向量空间，横截比较映射为复线性同构，其实行列式为复行列式绝对值的平方，严格为正。因此每个点的符号为 \(+1\)。最后，光滑旗模型紧，横截性在参数的小变动下保持；同时交点避开各闭包的低维部分，这一避开条件也由紧性保持。故所选好参数构成开集，参数横截性给出的稠密性则使它成为开稠密集。
\end{proof}

为了从这些计数读出一个类的系数，需要知道哪两个 Schubert 类相互对偶。定义矩形中的补分拆
\[
	\lambda^c=(q-\lambda_k,q-\lambda_{k-1},\ldots,q-\lambda_1).
\]
它把矩形中余下的方格旋转半周再左对齐，满足 \(|\lambda|+|\lambda^c|=N\)。

\begin{proposition}\label{g2:prop:schubert-complement-pairing}
	若 \(|\lambda|+|\mu|=N\)，则
	\begin{equation}\label{g2:eq:schubert-complement-pairing}
		\langle\sigma_\lambda\smile\sigma_\mu,[G]\rangle
		=\begin{cases}1,&\mu=\lambda^c,\\0,&\mu\ne\lambda^c.\end{cases}
	\end{equation}
\end{proposition}
\begin{proof}
	采用相反的两条坐标旗
	\[
		F_d=\operatorname{span}(e_1,\ldots,e_d),\qquad
		F'_d=\operatorname{span}(e_{n-d+1},\ldots,e_n).
	\]
	若某个 \(i\) 满足 \(\lambda_i+\mu_{k+1-i}>q\)，则相应两个旗子空间的维数和不超过 \(n\)，且它们在相反坐标旗中不相交。然而若 \(V\) 同时满足两组秩条件，它在这两个不相交子空间中分别至少有 \(i\) 维和 \(k+1-i\) 维，总计至少 \(k+1\) 维，矛盾。因此交集非空要求全部 \(\lambda_i+\mu_{k+1-i}\le q\)。将这些不等式求和，左边为 \(N=kq\)，所以每一项都必须取等号，即 \(\mu=\lambda^c\)。

	在补分拆情形，记 \(d_i=q+i-\lambda_i\)。上面的两子空间分别是 \(F_{d_i}\) 与 \(\operatorname{span}(e_{d_i},\ldots,e_n)\)，交集恰为 \(\mathbb Ce_{d_i}\)。由于 \(V\) 在二者中的维数之和至少为 \(k+1\)，它必须包含 \(e_{d_i}\)。对全部 \(i\) 使用此论证，唯一可能的交点便是
	\[
		V_0=\operatorname{span}(e_{d_1},\ldots,e_{d_k}).
	\]
	它确实满足两组条件，并位于两个开胞腔中。

	最后核对横截性。取补空间由非枢轴坐标向量张成，图坐标中的切向量为 \(A:V_0\to\mathbb C^n/V_0\)。由式\eqref{g2:eq:schubert-echelon-coordinates}，第一片的切方向是 \(A(e_{d_i})\) 在非枢轴行 \(a<d_i\) 的系数；相反旗的第二片恰好允许非枢轴行 \(a>d_i\) 的系数。这两组参数不交且合起来覆盖全部图坐标，所以切空间互补。复定向使唯一交点的符号为正，由前一命题的相交解释得到式\eqref{g2:eq:schubert-complement-pairing}。
\end{proof}

因此，若某个齐次类写成 \(\alpha=\sum_\lambda a_\lambda\sigma_\lambda\)，则
\[
	a_\lambda=\langle\alpha\smile\sigma_{\lambda^c},[G]\rangle.
\]
一个乘积的系数就能通过增加一项互补几何条件来求出。下面在最小的非射影空间例子中把这些条件逐项算完。

\subsection{\texorpdfstring{\(\operatorname{Gr}(2,4)\) 的杯积计算与四条直线问题}{Gr(2,4) 的杯积计算与四条直线问题}}
\label{g2:subsec:2105:05}
把 \(\operatorname{Gr}(2,4)\) 的点 \(V\) 看成 \(\mathbb{CP}^3\) 中的射影直线 \(\mathbb P(V)\)。固定射影旗“点 \(P\) 位于直线 \(L\)，而 \(L\) 位于平面 \(H\)”后，六个分拆的含义如下。
\begin{center}
	\begin{tabular}{ccl}
		\hline
		分拆 & 复余维 & 所参数化的射影直线 \\
		\hline
		\(\varnothing\) & \(0\) & 全部直线 \\
		\((1)\) & \(1\) & 与 \(L\) 相交 \\
		\((2)\) & \(2\) & 经过 \(P\) \\
		\((1,1)\) & \(2\) & 包含于 \(H\) \\
		\((2,1)\) & \(3\) & 经过 \(P\) 且包含于 \(H\) \\
		\((2,2)\) & \(4\) & 等于给定直线 \(L\) \\
		\hline
	\end{tabular}
\end{center}
其中 \(\Omega_2\) 与 \(\Omega_{11}\) 都同构于 \(\mathbb{CP}^2\)，\(\Omega_{21}\) 是一束直线组成的 \(\mathbb{CP}^1\)，而 \(\Omega_{22}\) 是一点。\(\Omega_1\) 的闭包却不必光滑，因而前面的光滑代表与移动旗论证在此已有实际用途。为简化式子，省略下标分拆中的逗号。

\begin{proposition}\label{g2:prop:gr24-schubert-products}
	\(G=\operatorname{Gr}(2,4)\) 的非零正次乘积由以下等式及交换律决定：
	\begin{equation}\label{g2:eq:gr24-schubert-products}
	\begin{gathered}
		\sigma_1^2=\sigma_2+\sigma_{11},\qquad
		\sigma_1\sigma_2=\sigma_1\sigma_{11}=\sigma_{21},\\
		\sigma_1\sigma_{21}=\sigma_{22},\qquad
		\sigma_2^2=\sigma_{11}^2=\sigma_{22},\qquad
		\sigma_2\sigma_{11}=0.
	\end{gathered}
	\end{equation}
	特别地，\(\sigma_1^3=2\sigma_{21}\)、\(\sigma_1^4=2\sigma_{22}\)，且 \(\langle\sigma_{22},[G]\rangle=1\)。
\end{proposition}
\begin{proof}
	先求余维互补的乘积。经过两个不同点的射影直线唯一，所以 \(\sigma_2^2=\sigma_{22}\)。同时位于两个不同平面的直线是它们的交线，所以 \(\sigma_{11}^2=\sigma_{22}\)。点若不在给定平面中，就没有既经过该点又包含于该平面的直线，故 \(\sigma_2\sigma_{11}=0\)。这也与补分拆配对相符：\((2)\) 与 \((1,1)\) 各自为自身的补。

	在平面 \(H\) 中经过点 \(P\) 的直线组成一束。另取一条不包含于 \(H\) 的直线 \(L\)，并使交点 \(Q=L\cap H\) 不等于 \(P\)。这束直线中恰有连接 \(P,Q\) 的一条与 \(L\) 相交，所以 \(\sigma_1\sigma_{21}=\sigma_{22}\)。一般位置保证交点横截，复定向使系数为一。

	六次群只由 \(\sigma_{21}\) 生成。若 \(P\notin L\)，经过 \(P\) 且与 \(L\) 相交的直线组成平面 \(\langle P,L\rangle\) 中的直线束，正是某条旗的 \(\Omega_{21}\)。在一般位置下相交横截于光滑部分，因而 \(\sigma_1\sigma_2=\sigma_{21}\)。同理，平面 \(H\) 内与一条横穿 \(H\) 的直线相交，等价于经过 \(L\cap H\)，所以 \(\sigma_1\sigma_{11}=\sigma_{21}\)。也可把这两个等式分别与 \(\sigma_1\) 配对，由刚才的唯一交点计数及非退化性确定系数。

	最后写 \(\sigma_1^2=a\sigma_2+b\sigma_{11}\)。分别乘 \(\sigma_2\) 与 \(\sigma_{11}\) 并在 \([G]\) 上求值，前面已经算得的乘积给出 \(a=b=1\)。其余两式由继续乘 \(\sigma_1\) 得到。次数超过八的类为零，因此没有其他正次乘积需要决定。
\end{proof}

这里的两种余维二条件不能混为一类：经过一个点与位于一个平面都给出一份 \(\mathbb{CP}^2\)，但它们在 \(G\) 中有不同同调类。它们各自的自交为一，彼此的交数为零，正是四次群秩为二的几何表现。

\ProseParagraphBreak

也可以将乘法表写成陈类的环表示。记 \(h=\sigma_1\)、\(s=\sigma_{11}\)。取一个非零线性函数 \(\alpha:\mathbb C^4\to\mathbb C\)，它限制到各 \(V\) 上给出 \(S^*\) 的截面。零点集是 \(V\subset\ker\alpha\)，即 \(\Omega_{11}\)。在这个零集处，图坐标微分为 \(A\mapsto\alpha\circ A\)，满射到 \(V^*\)，所以截面横截；于是
\[
	s=c_2(S^*),\qquad c(S)=1-h+s,
	\qquad c(Q)=1+h+(h^2-s).
\]
其中 \(c_2(Q)=\sigma_2=h^2-s\)，与特殊 Schubert 类公式一致。乘法表等价于
\begin{equation}\label{g2:eq:gr24-cohomology-ring}
	H^*(\operatorname{Gr}(2,4);\mathbb Z)
	\cong\mathbb Z[h,s]/(h^3-2hs,\ s^2-h^2s),
	\qquad\deg h=2,\quad\deg s=4.
\end{equation}
为核对没有漏掉整数关系，先用两条关系推出 \(h^3s=0\)、\(s^3=0\)、\(h^5=0\)：例如 \(h^3s=2hs^2=2h^3s\)，相减即得第一式。故商环由 \(1,h,h^2,s,hs,s^2\) 加法生成；这些元素的像分别为 \(1,\sigma_1,\sigma_2+\sigma_{11},\sigma_{11},\sigma_{21},\sigma_{22}\)，构成同调所给自由基的另一组整数基，所以映射既满又单。

\begin{theorem}\label{g2:thm:four-lines-two-transversals}
	在 \(\mathbb{CP}^3\) 中，对处于一般位置的四条射影直线，恰有两条射影直线同时与它们相交；两条解不同，每条解的局部相交重数为一。
\end{theorem}
\begin{proof}
	每条给定直线产生一份 \(\Omega_1\) 条件。由\cref{g2:prop:schubert-moving-flags}，一般位置下四份条件横截相交于各自开胞腔，故都是孤立的光滑交点，局部符号为 \(+1\)。总数于是为
	\[
		\langle\sigma_1^4,[G]\rangle
		=2\langle\sigma_{22},[G]\rangle=2.
	\]
	横截性的局部方程具有可逆复 Jacobian 矩阵，故每个零点都是单零点，其局部重数为一。
\end{proof}

这个数字还可以在一个直观的二次曲面模型中看出。先取三条两两不交的直线。前两条来自分解 \(\mathbb C^4=A\oplus B\) 的两个二维因子；第三条由于与它们都不交，可以写成某个同构 \(T:A\to B\) 的图。借助 \(T\) 识别 \(B\cong A\)，三条线便为
\[
	L_1=\mathbb P(A\oplus0),\quad L_2=\mathbb P(0\oplus A),\quad
	L_3=\{[a:a]:[a]\in\mathbb P(A)\}.
\]
一条同时遇到前两条线的直线必连接 \([a:0]\) 与 \([0:b]\)。它再与第三条相交，当且仅当 \(a,b\) 成比例。因此三条线的公共截线恰为
\[
	\ell_{[a]}=\{[sa:ta]:[s:t]\in\mathbb{CP}^1\},\qquad[a]\in\mathbb P(A).
\]
这些直线扫出二次曲面
\begin{equation}\label{g2:eq:three-lines-ruled-quadric}
	Q=\{[x_1:x_2:y_1:y_2]:x_1y_2-x_2y_1=0\}
	\cong\mathbb{CP}^1\times\mathbb{CP}^1.
\end{equation}
这个同构就是 \(([a],[s:t])\mapsto[sa:ta]\)：秩一的二乘二矩阵可分解为列向量乘行向量，两个因子分别只差非零倍数，所以参数唯一。非零秩一矩阵处方程的梯度不为零，故曲面光滑。固定 \([s:t]\) 得到包含原三条线的一个线族，固定 \([a]\) 得到公共截线的另一个线族。

\ProseParagraphBreak

一般的第四条线不包含于 \(Q\)，其参数代入二次方程后得到一个有两个不同根的齐次二次式。因而它与 \(Q\) 相交于两个点，每点唯一确定一条 \(\ell_{[a]}\)。这两条截线不同：若两个点属于同一条 \(\ell_{[a]}\)，则第四条线就是这条截线，因而包含于 \(Q\)，与选择矛盾。这给出了两条解的具体几何位置。若第四条线与 \(Q\) 相切，两个交点可能合并；若第四条线本身属于适当线族，还可能出现整族解。一般位置在这里排除的正是这些退化。

\subsection{与 V.27.5 中 Chow 环及 Plücker 方程计算的比较}
\label{g2:subsec:2105:06}
Grassmann 流形也有一个具体的射影方程模型。定义 Pl\"ucker 映射\footnote{Julius Pl\"ucker（1801--1868），德国数学家，研究解析几何、直线几何及数学物理。}
\[
	\iota:G\longrightarrow\mathbb P(\textstyle\bigwedge^k\mathbb C^n),
	\qquad\operatorname{span}(v_1,\ldots,v_k)\longmapsto[v_1\wedge\cdots\wedge v_k].
\]
换基只将外积乘以行列式，所以射影点不变。非零可分解外积 \(\xi=v_1\wedge\cdots\wedge v_k\) 又能恢复原子空间为 \(\{v:v\wedge\xi=0\}\)，故映射单射。在一个非零 \(k\) 阶子式对应的图域中，图矩阵的各项都是其他子式与这一子式之比，所以逆坐标光滑且微分单射。\(G\) 紧，因而 \(\iota\) 是到其像的光滑复嵌入。

\ProseParagraphBreak

其拉回的射影重言线在 \(V\) 处正是 \(\bigwedge^kV\)，所以
\[
	\iota^*\mathcal O(1)=\det S^*,\qquad\iota^*c_1(\mathcal O(1))=\sigma_1.
\]
这里 \(\mathcal O(1)\) 表示复射影空间的超平面线丛。若固定一个 \(q\) 维子空间 \(A\)，投影 \(\mathbb C^n\to\mathbb C^n/A\) 的行列式在 \(\bigwedge^k\mathbb C^n\) 上给出线性函数。它在 \(V\) 的 Pl\"ucker 点处为零，恰好当 \(V\cap A\ne0\)，所以特殊除子 \(\Omega_1\) 确为这一嵌入下的超平面截面。杯积中的每个 \(\sigma_1\)，因此也可以理解为加入一条这样的线性方程。

\ProseParagraphBreak

对 \(\operatorname{Gr}(2,4)\)，令 \(p_{ij}\) 为外积的六个坐标。由 \(\xi\wedge\xi=0\) 得到
\begin{equation}\label{g2:eq:gr24-plucker-quadric}
	p_{12}p_{34}-p_{13}p_{24}+p_{14}p_{23}=0.
\end{equation}
反过来，若某个坐标如 \(p_{12}\ne0\)，把它归一化为一，便可取矩阵
\[
	\begin{pmatrix}
		1&0&-p_{23}&-p_{24}\\
		0&1&p_{13}&p_{14}
	\end{pmatrix}.
\]
它的前五个相应子式是指定坐标，最后一个为 \(p_{13}p_{24}-p_{14}p_{23}=p_{34}\)。其他非零坐标图只需重新排列基，所以式\eqref{g2:eq:gr24-plucker-quadric}正好刻画 Pl\"ucker 像。因而这个 Grassmann 流形是 \(\mathbb{CP}^5\) 中的光滑二次超曲面，\(\int_G\sigma_1^4=2\) 与其二次方程相呼应。

第五卷第27章第5节将沿这个同一模型组织 Chow 环与方程计算：子簇先按有理等价形成代数循环类，再通过复几何的循环类映射与这里的整数上同调比较，并处理具体方程的消去、局部重数和退化参数。当前已经建立的是胞腔基本类、特殊陈类与横截相交计数；一般位置下每个点的局部贡献已由复横截性确定为一。转到代数模型后，重数还要在点合并或交集不横截时继续记录总计数，两种描述因此共享同一算例而承担不同的计算任务。
% 纲要条目排为编号小节标题，并进入目录与书签；说明段仍保留于纲要注释中。


\begin{chaptersummary}
	实线丛的第一 Stiefel--Whitney 类记录方向沿闭路的变化，复线丛的第一陈类是其底层定向实平面丛的 Euler 类。对偶重言复线丛在复射影直线上的第一类取值为 \(1\)，这固定了本章的符号。过渡函数的局部辐角把第一陈类表示为整数二上闭链，实线丛复化则给出整数 Bockstein 类。

	\ProseParagraphBreak

	Leray--Hirsch 定理将整体纤维基转成总空间的自由上同调模。射影丛关系的系数定义各阶陈类与 Stiefel--Whitney 类；迭代射影化所得的旗丛使原丛分裂为线丛，且上同调拉回单射。自然性、Whitney 公式及根的对称计算因而成立。顶陈类等于复定向 Euler 类，顶 Stiefel--Whitney 类等于模二 Euler 类；Thom 平方得到同一组实示性类，对角嵌入又给出 \(w(\mathrm TM)=\operatorname{Sq}(v(M))\)。

	\ProseParagraphBreak

	复化定义整数 Pontryagin 类。它们稳定、自然，但一般直和的整数公式需要保留奇次陈类乘积造成的二阶扭修正。射影空间的切丛稳定同构使三类计算具体化；球面的稳定平凡性则说明这些稳定类与依赖原秩的 Euler 类各有作用。示性类还能给出附加结构的条件：\(w_1\) 检测方向，\(c_1\) 检测行列式平凡性，而已定向丛的 spin 提升恰由 \(w_2\) 控制。

	\ProseParagraphBreak

	复 Grassmann 流形的 Schubert 胞腔给出整数加性基，几何入射条件的对偶类提供乘法解释。特殊 Schubert 类等于商丛的陈类，因而线性子空间之间的相交问题可以转成示性类运算。在 \(\operatorname{Gr}_2(\mathbb C^4)\) 中，四次超平面型 Schubert 条件的乘积求值为 \(2\)，对应三维复射影空间中同时与四条一般位置直线相交的两条直线。
\end{chaptersummary}

\BookExercises
\begin{exercise}\label{g2:ex:line-sections-positive-negative-indices}
	取整数 \(p,q\ge0\)，\(p+q=d>0\)。在 \(\mathbb CP^1\) 上，齐次单项式 \(z_0^pz_1^q\) 给出 \(\mathcal O(d)\) 的一个截面。找出全部零点，计算局部指标，并判断这些零点何时横截。再构造 \(\mathcal O(-d)\) 的光滑截面，使零点仍位于同样的位置，且非零局部指标全部为负。
\end{exercise}
\begin{solution}
	\(\mathcal O(d)=(\gamma^*)^{\otimes d}\)。在每条复直线 \(\ell\subset\mathbb C^2\) 上，齐次多项式限制为一个齐次 \(d\) 次函数，即 \((\ell^*)^{\otimes d}\) 的元素，所以题给式确实定义一个整体截面。

	当 \(p>0\) 时，\([0:1]\) 是零点。在附近以 \(w=z_0/z_1\) 作坐标，以 \((w,1)\) 及其对偶的张量幂作标架，截面系数为 \(w^p\)。在小圆周上归一化后，它是 \(e^{\mathrm i\theta}\mapsto e^{\mathrm i p\theta}\)，故局部指标为 \(p\)。同理，当 \(q>0\) 时，\([1:0]\) 是零点，以 \(v=z_1/z_0\) 作坐标，局部系数为 \(v^q\)，指标为 \(q\)。指数为零的那一端不是零点，此外没有零点。因此指标和是 \(p+q=d\)，与 \(c_1(\mathcal O(d))=dh\) 的求值相同。

	局部映射 \(w\mapsto w^p\) 在零点的实微分可逆恰好当 \(p=1\)；若 \(p>1\)，微分为零。另一端同理。所以所有实际零点横截恰好当 \(p,q\) 中每个非零指数都为一。指标和计入重数，而非只数零点位置。

	将此截面作纤维复共轭，得到 \(\overline{\mathcal O(d)}\) 的光滑截面。Hermitian 内积给出复线性同构 \(\overline{\mathcal O(d)}\cong\mathcal O(d)^*=\mathcal O(-d)\)。在共轭标架中的局部系数分别为 \(\overline w^p\)、\(\overline v^q\)；其归一化圆周映射的度分别为 \(-p\)、\(-q\)。随后施加的丛同构在各纤维中保持复方向，故不改变这些指标。所得截面的指标和为 \(-d\)，零点位置仍与原截面相同。
\end{solution}

\begin{exercise}\label{g2:ex:circle-bundle-torsion-line}
	在 \(\mathbb CP^1\) 上取 \(\mathcal O(m)\)，其中 \(m\ge2\)，令 \(Y_m=S(\mathcal O(m))\) 为其单位圆丛，投影为 \(\pi\)。对 \(Y_m\) 上的线丛 \(L=\pi^*\mathcal O(1)\)，计算 \(c_1(L)\) 的阶，并确定哪些整数 \(k\) 使 \(L^{\otimes k}\) 平凡。说明 \(c_1(L)\) 换成实系数后为何为零，并直接写出 \(L^{\otimes m}\) 的非零截面。
\end{exercise}
\begin{solution}
	\(Y_m\) 是紧光滑流形，第17章的 Morse 胞腔模型使它具有有限 CW 同伦型；结合丛的同伦不变性，可以在它上面使用线丛的上同调分类。记 \(h=c_1(\mathcal O(1))\)。给底层实二维丛取复方向，其 Euler 类为 \(mh\)。\cref{g2:thm:thom-gysin}的相应一段为
	\[
		H^0(\mathbb CP^1;\mathbb Z)=\mathbb Z
		\xrightarrow{\ \times m\ }H^2(\mathbb CP^1;\mathbb Z)=\mathbb Z
		\xrightarrow{\pi^*}H^2(Y_m;\mathbb Z)\longrightarrow0.
	\]
	于是 \(H^2(Y_m;\mathbb Z)\cong\mathbb Z/m\)，\(c_1(L)=\pi^*h\) 是其生成元，恰有阶 \(m\)。线丛的张量分类给出
	\[
		L^{\otimes k}\text{ 平凡}\quad\Longleftrightarrow\quad k\,c_1(L)=0
		\quad\Longleftrightarrow\quad m\mid k.
	\]
	负整数 \(k\) 按对偶解释，结论不变。换成实系数后，乘 \(m\) 是 \(\mathbb R\) 上的同构，同一段序列给出 \(H^2(Y_m;\mathbb R)=0\)，所以该非零整数扭类的实像为零。

	还可以直接看出第 \(m\) 个张量幂平凡。由定义，\(L^{\otimes m}=\pi^*\mathcal O(m)\)。一个点 \(y\in Y_m\) 本身就是 \(\mathcal O(m)_{\pi(y)}\) 中的单位向量，因此
	\[
		s(y)=(y,y)\in\bigl(\pi^*\mathcal O(m)\bigr)_y
	\]
	是处处非零的连续截面。它给出完整线标架，从而直接产生这个张量幂的平凡化。
\end{solution}

\begin{exercise}\label{g2:ex:klein-line-complexification}
	给 Klein 瓶 \(K\) 取一个零胞腔、两个有向一胞腔 \(a,b\) 及一个按字 \(aba^{-1}b\) 附着的二胞腔。令 \(\alpha,\beta\in H^1(K;\mathbb F_2)\) 分别在 \(a,b\) 上取值为 \((1,0)\)、\((0,1)\)。对 \(\varepsilon,\eta\in\mathbb F_2\)，记实线丛 \(L_{\varepsilon,\eta}\) 满足 \(w_1(L_{\varepsilon,\eta})=\varepsilon\alpha+\eta\beta\)。计算这些实线丛复化后的第一陈类，判断复化何时平凡，并说明四个不同实线丛会得到多少种复线丛。
\end{exercise}
\begin{solution}
	附着字中 \(a\) 的代数次数为零，\(b\) 的代数次数为二，所以整数胞腔边界为
	\[
		\partial_2(1)=2b,\qquad\partial_1=0.
	\]
	因此整数胞腔上链的非零余边界是
	\[
		\mathrm d^1:\mathbb Z^2\longrightarrow\mathbb Z,\qquad(u,v)\longmapsto2v,
	\]
	从而 \(H^2(K;\mathbb Z)=\mathbb Z/2\)。记二胞腔对偶上链所代表的生成元为 \(z\)。模二余边界全为零，故 \(\alpha,\beta\) 确实构成一次群的一组基。

	按整数 Bockstein 的除二公式，将 \(\alpha\) 提升为整数上链 \((1,0)\)，其余边界为零，所以 \(\widetilde\beta_2(\alpha)=0\)。将 \(\beta\) 提升为 \((0,1)\)，余边界为 \(2\)，除以二后得到二次上链 \(1\)，所以 \(\widetilde\beta_2(\beta)=z\)。由\cref{g2:prop:real-line-complexification-bockstein}，
	\[
		c_1(L_{\varepsilon,\eta}\otimes_{\mathbb R}\mathbb C)=\eta z.
	\]
	因此恰在 \(\eta=0\) 时复化平凡。原来的四个实线丛由四个不同的 \(w_1\) 区分；复化后只有两种同构类，\(\varepsilon\) 的差别消失，\(\eta\) 的差别保留为整数二阶扭类。这个计算在同一底空间上同时展示了实线丛非平凡性被复化消去和被复化保留的两种情况。
\end{solution}

\begin{exercise}\label{g2:ex:leray-hirsch-orientation-torus}
	令 \(a:S^2\to S^2\) 为对径映射，并取映射环面
	\[
		Y=(S^2\times[0,1])/\bigl((x,1)\sim(a(x),0)\bigr),\qquad p:Y\to S^1.
	\]
	判断纤维的整数二次生成元能否延拓为 \(Y\) 上的整数类，计算 \(H^*(Y;\mathbb Z)\)。再证明模二系数下存在满足 Leray--Hirsch 条件的整体纤维基，并计算模二上同调环。
\end{exercise}
\begin{solution}
	对径映射在 \(H^2(S^2;\mathbb Z)\) 上作用为 \(-1\)，因为\cref{g2:prop:reflection-antipodal-degree}给出它的度为 \((-1)^3=-1\)。若整数生成元来自整体类，沿基圆一周运输后必须保持不变，这与生成元变号矛盾。

	用基圆的零胞腔及一胞腔计算局部系数，\(q=0\) 行是常值 \(\mathbb Z\)，\(q=2\) 行的上链微分是 \(-2:\mathbb Z\to\mathbb Z\)。因此 Serre 谱序列的非零第二页项只有
	\[
		E_2^{0,0}=\mathbb Z,\qquad E_2^{1,0}=\mathbb Z,\qquad E_2^{1,2}=\mathbb Z/2.
	\]
	底空间只有两个可能的列，全部高阶微分为零；每个总次数也至多有一个非零项，故没有群延拓选择。于是
	\[
		H^k(Y;\mathbb Z)=
		\begin{cases}
			\mathbb Z,&k=0,1,\\
			\mathbb Z/2,&k=3,\\
			0,&\text{其他次数}.
		\end{cases}
	\]
	特别地，\(H^2(Y;\mathbb Z)=0\) 也直接排除了整数纤维生成元的延拓。

	模二系数下，\(-1=1\)，两行都是常值 \(\mathbb F_2\)。同一谱序列仍无非零微分，且边缘映射
	\[
		H^2(Y;\mathbb F_2)\longrightarrow H^2(S^2;\mathbb F_2)
	\]
	为同构。选 \(u\) 映到纤维生成元，则 \(1,u\) 满足 Leray--Hirsch 条件。令 \(x=p^*\eta\)，其中 \(\eta\) 是基圆的一次非零类。于是 \(1,x,u,xu\) 为加性基。\(x^2=p^*(\eta^2)=0\)，而 \(u^2=0\)，因为四次群为零；\(xu\ne0\) 由 Leray--Hirsch 的唯一表达保证。因此
	\[
		H^*(Y;\mathbb F_2)\cong\mathbb F_2[x,u]/(x^2,u^2),\qquad |x|=1,\quad |u|=2.
	\]
	同一纤维丛在不同系数下是否具有整体纤维基，取决于运输作用是否保留所需生成元。
\end{solution}

\begin{exercise}\label{g2:ex:projectivization-line-twist}
	设 \(E\to B\) 为复秩 \(r\ge1\) 丛，\(L\to B\) 为复线丛。构造覆盖底空间恒等映射的自然同构 \(\mathbb P(E)\cong\mathbb P(E\otimes L)\)，并比较两张重言线丛及其对偶的一阶类。若 \(r=2\)，将两条射影丛关系的系数分别记为 \(a_1,a_2\) 和 \(a'_1,a'_2\)。用 \(\ell=c_1(L)\) 表示 \(a'_1,a'_2\)。
\end{exercise}
\begin{solution}
	在每个 \(b\in B\) 上，将一条线 \(V\subset E_b\) 送到 \(V\otimes L_b\subset E_b\otimes L_b\)。局部选取 \(L\) 的非零标架后，这是通常的线子空间识别；更换标架只乘一个非零复数，不改变线，所以各局部映射粘合。张量 \(L_b^*\) 给出连续逆，因此得到自然同构 \(f:\mathbb P(E)\to\mathbb P(E\otimes L)\)。按每个点的纤维定义，有
	\[
		f^*\lambda_{E\otimes L}\cong\lambda_E\otimes p^*L.
	\]
	令 \(h=c_1(\lambda_E^*)\)、\(h'=f^*c_1(\lambda_{E\otimes L}^*)\)。线丛的张量及对偶公式给出
	\[
		h'=h-p^*\ell.
	\]
	所以射影丛作为空间相同，其用于记录原丛的重言线类仍可能不同。

	当 \(r=2\) 时，原关系为 \(h^2+p^*a_1h+p^*a_2=0\)。代入 \(h=h'+p^*\ell\)，整理得
	\[
		(h')^2+p^*(a_1+2\ell)h'+p^*(a_2+a_1\ell+\ell^2)=0.
	\]
	由射影关系系数的唯一性，
	\[
		a'_1=a_1+2\ell,\qquad a'_2=a_2+a_1\ell+\ell^2.
	\]
	这也说明原向量丛的信息不仅在射影总空间中，还在其重言线丛及投影中。
\end{solution}

\begin{exercise}\label{g2:ex:exterior-square-rank-three-roots}
	设 \(E\to B\) 为复秩三向量丛，记本章构造的陈类为 \(c_1,c_2,c_3\)。用分裂原理计算 \(\bigwedge^2E\) 的三个陈类，说明计算中的每个整数系数如何出现。
\end{exercise}
\begin{solution}
	在使 \(E\) 分裂的旗空间上写 \(q^*E=L_1\oplus L_2\oplus L_3\)，令 \(x_i=c_1(L_i)\)。外积的双线性给出实际丛分解
	\[
		\bigwedge^2(q^*E)\cong(L_1\otimes L_2)\oplus(L_1\otimes L_3)\oplus(L_2\otimes L_3).
	\]
	线丛张量的一阶类可加，所以新丛的三个形式根是
	\[
		y_1=x_1+x_2,\qquad y_2=x_1+x_3,\qquad y_3=x_2+x_3.
	\]
	它们的一次和为 \(2(x_1+x_2+x_3)\)，因为每个 \(x_i\) 出现两次。两两乘积之和为
	\[
	\begin{aligned}
		y_1y_2+y_1y_3+y_2y_3
		&=x_1^2+x_2^2+x_3^2+3(x_1x_2+x_1x_3+x_2x_3)\\
		&=(x_1+x_2+x_3)^2+(x_1x_2+x_1x_3+x_2x_3).
	\end{aligned}
	\]
	每个平方出现一次，每个混合项出现三次，所以最后恰剩一个二阶初等对称项。三个根的乘积为
	\[
		(x_1+x_2)(x_1+x_3)(x_2+x_3)
		=(x_1+x_2+x_3)(x_1x_2+x_1x_3+x_2x_3)-x_1x_2x_3.
	\]
	右边第一项中三重乘积出现三次，减去一次后留下左边的两次；其余六个三次单项式各出现一次。把初等对称多项式识别为 \(q^*c_i\)，再由 \(q^*\) 单射返回底空间，得到
	\[
		c_1\!\left(\bigwedge^2E\right)=2c_1,\qquad
		c_2\!\left(\bigwedge^2E\right)=c_1^2+c_2,\qquad
		c_3\!\left(\bigwedge^2E\right)=c_1c_2-c_3.
	\]
\end{solution}

\begin{exercise}\label{g2:ex:endomorphism-chern-low-classes}
	设 \(E\to B\) 为 CW 复形上的复秩 \(r\ge1\) 向量丛。计算复向量丛 \(\operatorname{End}(E)=E^*\otimes E\) 的第一、第二陈类，用 \(c_1(E),c_2(E)\) 表示结果。再说明将 \(\operatorname{End}(E)\) 换成无迹自同态丛 \(\operatorname{End}_0(E)\) 时，所有正次陈类为何不变。
\end{exercise}
\begin{solution}
	将 \(E\) 拉回到分裂空间，记其形式根为 \(x_1,\ldots,x_r\)。由\cref{g2:prop:chern-operations}，\(\operatorname{End}(E)\) 的根为全部 \(x_j-x_i\)。当 \(i=j\) 时根为零；其余根可按 \((i,j)\) 与 \((j,i)\) 配对，因此
	\[
		c(\operatorname{End}(E))
		=\prod_{i<j}\bigl(1+(x_j-x_i)\bigr)\bigl(1-(x_j-x_i)\bigr)
		=\prod_{i<j}\bigl(1-(x_j-x_i)^2\bigr)
	\]
	成立于分裂空间。每个因子的第一个正次项都是四次，所以 \(c_1(\operatorname{End}(E))=0\)。四次项给出
	\[
	\begin{aligned}
		c_2(\operatorname{End}(E))
		&=-\sum_{i<j}(x_j-x_i)^2\\
		&=-(r-1)\sum_i x_i^2+2\sum_{i<j}x_ix_j\\
		&=-(r-1)\bigl(c_1(E)^2-2c_2(E)\bigr)+2c_2(E)\\
		&=2r\,c_2(E)-(r-1)c_1(E)^2.
	\end{aligned}
	\]
	第二行中，每个 \(x_i^2\) 在含有下标 \(i\) 的 \(r-1\) 对中出现；混合项来自平方展开的 \(-2x_ix_j\)。旗丛拉回单射使这些等式返回原基空间。配对计算始终使用整数多项式恒等式，没有对上同调类除以二，故包含整数扭信息的基空间也适用。

	迹映射在换架共轭下不变，故 \(\operatorname{End}_0(E)=\ker\operatorname{tr}\) 是复秩 \(r^2-1\) 子丛。映射
	\[
		A\longmapsto\left(\frac{\operatorname{tr}A}{r},\ A-\frac{\operatorname{tr}A}{r}I\right)
	\]
	给出 \(\operatorname{End}(E)\cong\underline{\mathbb C}\oplus\operatorname{End}_0(E)\)，逆映射为 \((z,A_0)\mapsto zI+A_0\)。这里的除以 \(r\) 发生在复纤维的线性代数中。平凡线因子的总陈类为一，所以 Whitney 公式说明两张丛的全部正次陈类相同。
\end{solution}

\begin{exercise}\label{g2:ex:three-real-lines-projective-product}
	设 \(m,n\ge1\)，取 \(B=\mathbb RP^m\times\mathbb RP^n\)。记从两因子拉回的典范实线丛为 \(L_a,L_b\)，令 \(a=w_1(L_a)\)、\(b=w_1(L_b)\)，并置
	\[
		E=L_a\oplus L_b\oplus(L_a\otimes L_b).
	\]
	计算 \(w_1(E),w_2(E),w_3(E)\)，判断 \(E\) 是否可定向、是否平凡，并确定哪些 \((m,n)\) 使它存在处处非零截面。
\end{exercise}
\begin{solution}
	射影空间模二环及域系数 Künneth 公式给出
	\[
		H^*(B;\mathbb F_2)=\mathbb F_2[a,b]/(a^{m+1},b^{n+1}).
	\]
	第三条线的第一类为 \(a+b\)，所以 Whitney 公式得到
	\[
	\begin{aligned}
		w(E)&=(1+a)(1+b)(1+a+b),\\
		w_1(E)&=0,\\
		w_2(E)&=a^2+ab+b^2,\\
		w_3(E)&=a^2b+ab^2.
	\end{aligned}
	\]
	由\cref{g2:prop:first-classes-determinants}，\(w_1=0\) 说明 \(E\) 可定向。由于 \(m,n\ge1\)，单项式 \(ab\) 在上述截断多项式环中始终非零，并且与尚未消失的 \(a^2,b^2\) 线性无关。因此 \(w_2(E)\ne0\)，\(E\) 对所有给定 \((m,n)\) 都不平凡。

	当至少一个 \(m,n\) 大于一时，\(a^2b\)、\(ab^2\) 中至少有一个非零；若两项均非零，它们是不同的基单项式，不能相消。所以此时 \(w_3(E)\ne0\)。\cref{g2:thm:top-class-euler}把它识别为模二 Euler 类，而非零截面会使该类为零，故这些情形没有处处非零截面。

	剩下 \(m=n=1\) 时，底空间是二维环面，\(w_3(E)=0\)。给 \(E\) 选内积，单位球丛的纤维为 \(S^2\)。先在零骨架选单位向量；纤维道路连通，使该选择能沿每条一胞腔延拓。每张二维胞腔上的丛平凡，边界截面对应一条 \(S^1\to S^2\) 的映射，而 \(\pi_1(S^2)=0\)，故还能延入圆盘。底空间没有更高胞腔，得到整体单位截面。因此恰在 \((m,n)=(1,1)\) 时存在非零截面。此时 \(w_2(E)=ab\ne0\)，所以一条非零截面仍不足以形成整个三维丛的标架。
\end{solution}

\begin{exercise}\label{g2:ex:real-projective-wu-formula-all-dimensions}
	设 \(n\ge1\)，令 \(a\in H^1(\mathbb RP^n;\mathbb F_2)\) 为非零类。直接由 Wu 类的求值定义与射影空间上的平方公式证明
	\[
		v_i(\mathbb RP^n)=\binom{n-i}{i}a^i\qquad(0\le i\le n),
	\]
	其中二项式系数取模二，且下标大于非负上标时约定为零。说明 \(i>n/2\) 时为何消失，并分别计算 \(n=4,5\) 时的总 Wu 类及由 Wu 公式得到的切丛总 Stiefel--Whitney 类。
\end{exercise}
\begin{solution}
	对 \(0\le i\le n\)，\(H^i(\mathbb RP^n;\mathbb F_2)\) 由 \(a^i\) 生成，因而可写 \(v_i=\lambda_i a^i\)。模二基本类满足 \(\langle a^n,[\mathbb RP^n]_2\rangle=1\)，所以把 Wu 类的定义用于 \(a^{n-i}\)，得到
	\[
	\begin{aligned}
		\lambda_i
		&=\langle v_i a^{n-i},[\mathbb RP^n]_2\rangle\\
		&=\langle\operatorname{Sq}^i(a^{n-i}),[\mathbb RP^n]_2\rangle
		=\binom{n-i}{i}.
	\end{aligned}
	\]
	最后一步使用式\eqref{g2:eq:projective-steenrod-formula}。若 \(i>n/2\)，则 \(i>n-i\)，平方次数超过输入类的次数，不稳定条件使它为零；这也与二项式系数的约定一致。\(i>n\) 时所在上同调群本来为零。

	当 \(n=4\) 时，可能的正次项只有一次和二次，系数分别为 \(\binom31=1\)、\(\binom22=1\)（模二），故
	\[
		v(\mathbb RP^4)=1+a+a^2.
	\]
	由平方公式，\(\operatorname{Sq}(a)=a+a^2\)，\(\operatorname{Sq}(a^2)=a^2+a^4\)，中间的 \(\operatorname{Sq}^1(a^2)=2a^3=0\)。因此\cref{g2:thm:wu-formula}给出
	\[
		w(\mathrm T\mathbb RP^4)=\operatorname{Sq}(1+a+a^2)=1+a+a^4.
	\]
	当 \(n=5\) 时，\(\binom41=0\)、\(\binom32=1\)（模二），所以
	\[
		v(\mathbb RP^5)=1+a^2,\qquad
		w(\mathrm T\mathbb RP^5)=1+a^2+a^4.
	\]
	两种情形中 Wu 类的高次部分都先消失，随后对低次 Wu 类施加平方，仍能产生切丛较高次的非零类。两组结果也分别等于截断后的 \((1+a)^5\)、\((1+a)^6\)，与射影空间切丛的直接计算一致。
\end{solution}

\begin{exercise}\label{g2:ex:projective-stiefel-binary}
	对 \(n\ge1\)，证明 \(\mathbb RP^n\) 的全部正次 Stiefel--Whitney 类为零，当且仅当 \(n+1\) 是二的幂。再对 \(\mathbb RP^3\)、\(\mathbb RP^5\)、\(\mathbb RP^7\) 分别判断可定向性与 spin 结构的存在性，并计算存在时 spin 结构的个数。
\end{exercise}
\begin{solution}
	令 \(m=n+1\)，在 \(\mathbb F_2[a]/(a^m)\) 中有 \(w(\mathrm T\mathbb RP^n)=(1+a)^m\)。将 \(m\) 写成二进制 \(m=\sum_{j\in J}2^j\)。因为特征二下 \((1+a)^{2^j}=1+a^{2^j}\)，所以
	\[
		(1+a)^m=\prod_{j\in J}(1+a^{2^j}).
	\]
	若 \(m=2^k\)，右端为 \(1+a^m=1\)。反过来，若 \(J\) 至少有两个元素，令 \(j_0=\min J\)。右端 \(a^{2^{j_0}}\) 的系数为一：只能从第 \(j_0\) 个因子选择正次项，其他正次项的次数都更大。而 \(2^{j_0}<m\)，这一项在商环中不为零。因此总类不等于一，得到所求判据。

	三个维数都是奇数，所以相应实射影空间都可定向。\(w_2\) 的系数依次是 \(\binom42=6\)、\(\binom62=15\)、\(\binom82=28\)，故 \(\mathbb RP^3\)、\(\mathbb RP^7\) 有 spin 结构，\(\mathbb RP^5\) 没有。前两者的一次模二上同调均为 \(\mathbb F_2\)，所以各有两份 spin 结构。正次 Stiefel--Whitney 类全零是示性类层面的结论；判定实际切丛是否平凡，还需研究完整的标架延拓。
\end{solution}

\begin{exercise}\label{g2:ex:product-projective-chern-numbers}
	给 \(M=\mathbb CP^1\times\mathbb CP^2\) 乘积复结构。记两个因子的正生成元拉回为 \(x,y\)，所以 \(x^2=0\)、\(y^3=0\)，且 \(\langle xy^2,[M]\rangle=1\)。求 \(c_1,c_2,c_3\)、\(p_1(\mathrm TM)\)，计算三个陈数 \(\langle c_1^3,[M]\rangle\)、\(\langle c_1c_2,[M]\rangle\)、\(\langle c_3,[M]\rangle\)，并判断 \(M\) 是否有 spin 结构。
\end{exercise}
\begin{solution}
	乘积切丛分裂为两个因子的切丛拉回，所以
	\[
		c(\mathrm TM)=(1+x)^2(1+y)^3
		=(1+2x)(1+3y+3y^2).
	\]
	按次数收集得到
	\[
		c_1=2x+3y,\qquad c_2=6xy+3y^2,\qquad c_3=6xy^2.
	\]
	因此
	\[
		p_1=c_1^2-2c_2=(12xy+9y^2)-(12xy+6y^2)=3y^2.
	\]
	只有 \(xy^2\) 项在顶次求值中有贡献，于是
	\[
		c_1^3=54xy^2,\qquad c_1c_2=24xy^2,\qquad c_3=6xy^2.
	\]
	三个数分别为 \(54,24,6\)。最后 \(w_2(\mathrm TM)=\rho_2c_1=\overline y\ne0\)，其非零性可通过限制到第二个因子的射影直线检验。因此 \(M\) 没有 spin 结构。
\end{solution}

\begin{exercise}\label{g2:ex:two-spin-lifts-circle}
	设 \(E=S^1\times\mathbb R^3\) 带标准方向与度量。用 \(\operatorname{Spin}(3)=S^3\) 及四元数覆盖 \(\rho:S^3\to\operatorname{SO}(3)\)，对 \(\epsilon\in\{1,-1\}\) 定义
	\[
		Q_\epsilon=(\mathbb R\times S^3)/\!\sim,
		\qquad(t+1,q)\sim(t,\epsilon q).
	\]
	证明 \(\eta_\epsilon([t,q])=(e^{2\pi\mathrm i t},\rho(q))\) 给出 \(E\) 的 spin 结构。证明二者作为覆盖固定标架丛的 spin 结构不同构，却作为抽象主 \(S^3\) 丛都平凡，并说明二者为何穷尽所有 spin 结构。
\end{exercise}
\begin{solution}
	\(\epsilon\) 位于 \(S^3\) 的中心，因此右乘作用 \([t,q]g=[t,qg]\) 有定义，短弧上的实参数提升给出主丛局部平凡化。又 \(\rho(\epsilon q)=\rho(q)\)，故 \(\eta_\epsilon\) 尊重商关系，且逐纤维为 \(\rho\)、满足所需等变性。

	取固定标架丛中的截面 \(s(z)=(z,I_3)\)。其在 \(Q_\epsilon\) 中的原像是
	\[
		(\mathbb R\times\{1,-1\})/\bigl((t+1,a)\sim(t,\epsilon a)\bigr).
	\]
	\(\epsilon=1\) 时，这是两个互不相交的圆；\(\epsilon=-1\) 时，两张局部覆盖沿一周互换，得到一个连通的二重圆周覆盖。任何与 \(\eta\) 交换的同构都须把这两份原像同构地对应，连通分支数不同使它不可能存在。

	作为抽象主丛，\(Q_1\) 已经是乘积。对 \(Q_{-1}\)，取单位纯虚四元数 \(u\)，令 \(h(t)=\cos(\pi t)+u\sin(\pi t)\)，则 \(h(t+1)=-h(t)\)。映射
	\[
		[t,q]\longmapsto(e^{2\pi\mathrm i t},h(t)q)
	\]
	尊重接缝、保持右作用，并在局部具有连续逆，所以给出主丛平凡化。它改变了到固定 \(\operatorname{SO}(3)\) 标架丛的映射，因而不与前面不同构的结论冲突。由\cref{g2:thm:spin-obstruction}，所有 spin 结构同构类受 \(H^1(S^1;\mathbb F_2)=\mathbb F_2\) 自由传递作用，恰有两份，故上述构造已经穷尽。
\end{solution}

\begin{exercise}\label{g2:ex:smooth-projective-surface-numbers}
	设 \(F(z_0,z_1,z_2,z_3)\) 为次数 \(d\ge1\) 的复齐次多项式，其射影零点集 \(X\subset\mathbb CP^3\) 是光滑复曲面；等价地，所定义的齐次方程沿零点集的法向微分处处满射。记 \(h=c_1(\lambda^*)\)，并用同一字母表示其在 \(X\) 上的限制。证明
	\[
		c_1(\mathrm TX)=(4-d)h,\qquad
		c_2(\mathrm TX)=(d^2-4d+6)h^2,\qquad
		\langle h^2,[X]\rangle=d.
	\]
	据此计算 \(X\) 的两个陈数及第一 Pontryagin 数，并证明 \(X\) 有 spin 结构当且仅当 \(d\) 为偶数。
\end{exercise}
\begin{solution}
	在一条直线 \(\ell\subset\mathbb C^4\) 上，\(F\) 的限制是 \(\ell\) 上的次数 \(d\) 齐次函数，所以是 \((\ell^*)^{\otimes d}\) 的元素。因此 \(F\) 给出线丛 \((\lambda^*)^{\otimes d}\) 的光滑截面，零集正是 \(X\)。题设满秩条件说明该截面横截于零，并将复法丛识别为 \((\lambda^*)^{\otimes d}|_X\)。选 Hermitian 度量分裂切丛的短正合列，得到复丛同构
	\[
		\mathrm T\mathbb CP^3|_X\cong\mathrm TX\oplus(\lambda^*)^{\otimes d}|_X.
	\]
	线丛张量公式给出后一项的第一陈类为 \(dh\)，因此
	\[
		c(\mathrm TX)=\frac{(1+h)^4}{1+dh}
		=(1+4h+6h^2)(1-dh+d^2h^2)
		=1+(4-d)h+(d^2-4d+6)h^2,
	\]
	这里只保留实次数不超过四的项。

	由\cref{g2:thm:euler-transverse-zero-duality}，\(X\) 的对偶类是这条复线丛的 Euler 类，即 \(dh\)。复法向实维数为二，交换法向与切向的方向不会产生符号，所以该零点方向与 \(X\) 的复方向一致。于是
	\[
		\langle h^2,[X]\rangle=\langle h^2\smile dh,[\mathbb CP^3]\rangle=d.
	\]
	从而
	\[
	\begin{aligned}
		\langle c_1^2,[X]\rangle&=d(4-d)^2,\\
		\langle c_2,[X]\rangle&=d(d^2-4d+6),\\
		\langle p_1(\mathrm TX),[X]\rangle&=\langle c_1^2-2c_2,[X]\rangle=d(4-d^2).
	\end{aligned}
	\]
	最后 \(w_2(\mathrm TX)=(4-d)\overline h\)。\(d\) 偶数时它为零。\(d\) 奇数时，\(\langle\overline h^{\,2},[X]_2\rangle=1\)，故 \(\overline h\ne0\)，于是 \(w_2\ne0\)。由 spin 提升准则得到所求充要条件。比如 \(d=4\) 时两种陈数为 \(0,24\)，第一 Pontryagin 数为 \(-48\)；这些数均由上述切丛与零点类计算得到。
\end{solution}

\begin{exercise}\label{g2:ex:gr25-cells-pairing}
	列出 \(\operatorname{Gr}(2,5)\) 的各次整数 Betti 数与 Euler 示性数。给出四次、八次上同调的 Schubert 基，并写出两组基之间的相交配对矩阵。
\end{exercise}
\begin{solution}
	此时 \(k=2\)、\(q=3\)、\(N=6\)。矩形中的分拆按大小排列为
	\[
	\begin{array}{c|l}
		|\lambda|&\lambda\\\hline
		0&(0,0)\\
		1&(1,0)\\
		2&(2,0),(1,1)\\
		3&(3,0),(2,1)\\
		4&(3,1),(2,2)\\
		5&(3,2)\\
		6&(3,3).
	\end{array}
	\]
	故从零次到十二次的偶数 Betti 数依次为 \(1,1,2,2,2,1,1\)，全部奇数 Betti 数为零。Euler 示性数为这些数之和，即 \(10\)。这里每个分拆贡献一个偶维胞腔，所以也可直接数得 \(\binom52=10\)。

	四次上同调取基 \((\sigma_2,\sigma_{11})\)，八次取基 \((\sigma_{31},\sigma_{22})\)。补分拆公式为 \(\lambda^c=(3-\lambda_2,3-\lambda_1)\)，所以 \((2,0)^c=(3,1)\)、\((1,1)^c=(2,2)\)。按上述次序，相交配对矩阵为
	\[
		\begin{pmatrix}1&0\\0&1\end{pmatrix}.
	\]
	这里第二组基与第一组基所在次数不同，因为 Grassmann 流形实维数为十二，互补次数是四与八。
\end{solution}

\begin{exercise}\label{g2:ex:gr24-tangent-chern-numbers}
	在 \(G=\operatorname{Gr}(2,4)\) 上令 \(h=\sigma_1\)、\(s=\sigma_{11}\)。证明复切丛同构于 \(S^*\otimes Q\)，计算 \(c_1(\mathrm TG)\)、\(c_2(\mathrm TG)\)、\(c_4(\mathrm TG)\)，并将顶陈数与胞腔计算得到的 Euler 示性数比较。
\end{exercise}
\begin{solution}
	在 \(V\) 附近，另一个二维子空间可写成从 \(V\) 到补空间的线性映射的图。对这些图在零映射处求微分，得到 \(\mathrm T_VG=\operatorname{Hom}_{\mathbb C}(V,\mathbb C^4/V)\)。更换补空间不改变商值的一阶映射，所以这些识别粘成 \(\mathrm TG\cong S^*\otimes Q\)。

	在共同分裂空间上记 \(S^*\) 的根为 \(x_1,x_2\)，\(Q\) 的根为 \(y_1,y_2\)。由正文计算，
	\[
		x_1+x_2=y_1+y_2=h,\qquad x_1x_2=s,\qquad y_1y_2=h^2-s.
	\]
	切丛四个根为 \(x_i+y_j\)。它们之和为 \(4h\)，所以 \(c_1(\mathrm TG)=4h\)。平方和为
	\[
	\begin{aligned}
		\sum_{i,j}(x_i+y_j)^2
		&=2(h^2-2s)+2\bigl(h^2-2(h^2-s)\bigr)+2h^2\\
		&=2h^2.
	\end{aligned}
	\]
	由根的对称多项式恒等式 \(2c_2=c_1^2-\sum_{i,j}(x_i+y_j)^2\)，得到 \(2c_2=14h^2\)。整数上同调没有二扭，故 \(c_2(\mathrm TG)=7h^2\)。

	顶类是四根之积。先对 \(y_j\) 求积，并利用 \(x_i^2=hx_i-s\)，有
	\[
	\begin{aligned}
		c_4(\mathrm TG)
		&=\prod_{i=1}^2\bigl(x_i^2+hx_i+h^2-s\bigr)\\
		&=\prod_{i=1}^2(2hx_i+h^2-2s)\\
		&=3h^4-4h^2s+4s^2=3h^4.
	\end{aligned}
	\]
	最后一步使用 \(s^2=h^2s\)。共同旗丛拉回单射，因此这些根计算给出原空间上的等式。由 \(\int_Gh^4=2\)，顶陈数为 \(6\)。\(G\) 恰有六个偶维胞腔，Euler 示性数也为 \(6\)；两者相等是复切丛顶陈类、实 Euler 类与\cref{g2:cor:tangent-euler-characteristic}的具体核验。
\end{solution}

\begin{exercise}\label{g2:ex:schubert-divisor-singularity}
	令 \(A\subset\mathbb C^4\) 为固定二维子空间，\(\Omega_1=\{V\in\operatorname{Gr}(2,4):V\cap A\ne0\}\)。证明 \(\Omega_1\) 在点 \(V=A\) 处不光滑，并在该点附近给出它的方程。再说明正文中的关联空间 \(Z=\{(V,L):L\subset A\cap V\}\) 在这个点上的投影纤维是什么。
\end{exercise}
\begin{solution}
	选择补空间 \(B\)，把 \(A\) 附近的 \(V\) 写成线性映射 \(T:A\to B\) 的图。则 \(V\cap A\) 对应 \(\ker T\)，所以局部方程为
	\[
		\det T=t_{11}t_{22}-t_{12}t_{21}=0
	\]
	于 \(\operatorname{Hom}(A,B)\cong\mathbb C^4\) 中。在 \(T\ne0\) 且行列式为零时，\(T\) 的秩为一，行列式的至少一个偏导数非零，故该处为复三维光滑超曲面。

	在 \(T=0\) 处，集合是所有秩至多一矩阵组成的锥。若它在那里是复三维光滑子流形，则其切空间必须包含每条曲线 \(t\mapsto tE_{ij}\) 在零点的速度。四个矩阵单位 \(E_{ij}\) 都秩一，且线性张成全部 \(\mathbb C^4\)，这与切空间只有复维数三矛盾。因此原点确实是奇点，而不仅是这个方程的一阶导数恰好同时消失。

	投影到 \(V=A\) 时，条件为任意直线 \(L\subset A\)，所以纤维是 \(\mathbb P(A)\cong\mathbb{CP}^1\)。在其余 \(\Omega_1\) 的点上，\(V\cap A\) 一维，纤维只有一个点。关联空间把奇点处不唯一的交线选择保留下来；正文通过其基本类推前定义同一个 Schubert 类，因此闭包有奇点并不妨碍类的存在。
\end{solution}

\begin{exercise}\label{g2:ex:real-four-lines-count}
	在 \(\mathbb C^4=\mathbb C^2\oplus\mathbb C^2\) 中取射影直线
	\[
		L_1=\mathbb P(\mathbb C^2\oplus0),\quad L_2=\mathbb P(0\oplus\mathbb C^2),\quad L_3=\mathbb P\{(a,a):a\in\mathbb C^2\}.
	\]
	第四条直线取 \(L_T=\mathbb P\{(a,Ta):a\in\mathbb C^2\}\)。分别令
	\[
		T_+=\begin{pmatrix}0&2\\3&0\end{pmatrix},\qquad
		T_-=\begin{pmatrix}0&-1\\2&0\end{pmatrix}.
	\]
	证明每组四条线两两不交，并求所有复公共截线。两组数据都由实矩阵给出，它们各有多少条实公共截线？
\end{exercise}
\begin{solution}
	前面三条直线两两不交。图直线与 \(L_2\) 不交是因为第一分量为零会迫使 \(a=0\)；与 \(L_1\) 不交等价于 \(T\) 可逆，与 \(L_3\) 不交等价于 \(T-I\) 可逆。两个所给矩阵分别满足
	\[
		\det T_+=-6,\quad\det(T_+-I)=-5,\qquad
		\det T_-=2,\quad\det(T_--I)=3,
	\]
	所以四条线两两不交。

	前三条线的公共截线由正文给出，为 \(\ell_{[a]}=\{[sa:ta]:[s:t]\in\mathbb{CP}^1\}\)。它与 \(L_T\) 相交，恰好当 \(Ta\) 与 \(a\) 成比例，即 \([a]\) 是 \(T\) 的特征线。对 \(T_+\)，特征值为 \(\pm\sqrt6\)，可取特征向量 \(a_\pm=(2,\pm\sqrt6)\)。所以两条公共截线为 \(\ell_{[a_+]}\)、\(\ell_{[a_-]}\)，均由实向量张成，是两条不同的实射影直线。

	对 \(T_-\)，特征值为 \(\pm\mathrm i\sqrt2\)，对应特征向量可取 \((1,\mp\mathrm i\sqrt2)\)。仍有两条不同的复公共截线，它们互为复共轭，且都不由实二维子空间定义。因为实矩阵的实特征线必有实特征值，而这里没有实特征值，所以没有实公共截线。两个例子说明复交数 \(2\) 可以对应两条实解，也可以对应一对非实共轭解，不能直接把复计数当作实计数。
\end{solution}
